\exercisesection{Torsion Products}

\begin{enumerate}
\def\labelenumi{\arabic{enumi})}
\item
  We denote by B the ring A(ε) (X, p.~27); let C be the complex of (B, B)-bimodules such that \(C_{n} = B\) for all \(n \in \mathbf{Z}\) , \(d(x) = \varepsilon x\) for all \(x \in C\) . Show that the complex C has zero homology, but that H \(_{n}\) (C ⊗ \(_{B}\) C) is isomorphic to A for all \(n \in \mathbf{Z}\) . Deduce that one cannot remove the hypothesis “E bounded above” in lemma 1 and prop. 4, pp.~66-67.
\item
  Let \(\rho: A \to B\) be a ring homomorphism, M an A-module. Show that the following properties are equivalent:
\end{enumerate}

α) One has \(\operatorname{Tor}_1^{\mathbf{A}}(\rho_*(\mathbf{P}), \mathbf{M}) = 0\) for every right B-module P.

β) One has \(\operatorname{Tor}_1^{\mathbf{A}}(\rho_{*}(\mathbf{P}),\mathbf{M}) = 0\) for every monogenic right B-module P.

\(\gamma)\) The B-module \(\mathbf{B} \otimes_{\mathbf{A}} \mathbf{M}\) is flat, and one has \(\operatorname{Tor}_{1}^{\mathbf{A}}(\mathbf{B}, \mathbf{M}) = 0\) .

(To see that \(\gamma\) ) implies \(\alpha\) ), consider an exact sequence \(0 \to R \to L \to P \to 0\) of right B-modules, where L is free.)

\begin{enumerate}
\def\labelenumi{\arabic{enumi})}
\setcounter{enumi}{2}
\tightlist
\item
  Assume the ring A is commutative. Let I be a finite set and \((\mathbf{C}^{(i)}, d^{(i)})_{i \in I}\) a family of complexes of A-modules.
\end{enumerate}

\begin{enumerate}
\def\labelenumi{\alph{enumi})}
\tightlist
\item
  We equip the tensor product \(\otimes\) \(\mathbf{C}^{(i)}\) with its graduation of type \(\mathbf{Z}^1\) and with the A-endomorphisms
\end{enumerate}

\[
d _ {i} = \bigotimes_ {j \in I} d _ {i} ^ {(j)} \quad \text { with } \quad d _ {i} ^ {(j)} = 1 _ {\mathrm{C} ^ {(j)}} \quad \text { for } \quad j \neq i \quad \text { and } \quad d _ {i} ^ {(i)} = d ^ {(i)}.
\]

Show that one thus obtains an I-precomplex (X, exercise 13, p.~174); the associated I-complex is called the tensor product I-complex of the complexes \(\mathbf{C}^{(i)}\) . The complex associated with this I-complex is called the tensor product complex of the \(\mathbf{C}^{(i)}\) .

\begin{enumerate}
\def\labelenumi{\alph{enumi})}
\setcounter{enumi}{1}
\tightlist
\item
  Show that the choice of a total order on I allows one to define an isomorphism from the tensor product complex of the \(\mathbf{C}^{(i)}\) to the complex defined in X, p.~64.
\end{enumerate}

\begin{enumerate}
\def\labelenumi{\arabic{enumi})}
\setcounter{enumi}{3}
\tightlist
\item
  Let a be a two-sided ideal of A. Let M be an A-module, such that \(\operatorname{Tor}_1^A (\mathbf{A} / \mathfrak{a},\mathbf{M}) = 0\) and that the \((\mathbf{A} / \mathfrak{a})\) -module \(\mathbf{M} / \mathfrak{a}\mathbf{M}\) is projective.
\end{enumerate}

\begin{enumerate}
\def\labelenumi{\alph{enumi})}
\item
  Assume either that M is finitely presented and a is contained in the radical of A, or that a is nilpotent. Show that M is projective (use VIII, § 8, n° 3, cor. 3).
\item
  Assume that M is finitely generated, that one has \(\bigcap \mathfrak{a}^n = 0\) and that the (A/a)-module M/aM is free. Show
\end{enumerate}

that the A-module M is free.

\begin{enumerate}
\def\labelenumi{\alph{enumi})}
\setcounter{enumi}{2}
\tightlist
\item
  Give an example of a local ring A, of maximal ideal m, and of a finitely generated A-module M which is not flat such that \(\operatorname{Tor}_{1}^{A}(A/m, M) = 0\) .
\end{enumerate}

\begin{enumerate}
\def\labelenumi{\arabic{enumi})}
\setcounter{enumi}{4}
\tightlist
\item
  Assume there exists a (unitary) homomorphism of \(k\) -algebras \(\pi: \mathbf{A} \to k\) , which endows \(k\) with a structure of A-module. We consider the standard resolution \(\mathbf{B}(\mathbf{A}, k)\) (X, p.~58); for \(n \geqslant 0\) , we identify \(\mathbf{B}_n(\mathbf{A}, k)\) with \(\mathbf{A} \otimes_k \mathbf{A}^{\otimes^n}\) and we denote by \(a[a_1, ..., a_n]\) the element \(a \otimes a_1 \otimes \cdots \otimes a_n\) of \(\mathbf{A} \otimes_k \mathbf{A}^{\otimes^n}\) .
\end{enumerate}

\begin{enumerate}
\def\labelenumi{\alph{enumi})}
\tightlist
\item
  Show that the complex \(\mathbf{B}(\mathbf{A}, k) \otimes_k \mathbf{B}(\mathbf{A}, k)\) defines a resolution of the \((\mathbf{A} \otimes_k \mathbf{A})\) -module \(k\) . Show that the homomorphism \((\mathbf{A} \otimes_k \mathbf{A})\) -linear
\end{enumerate}

\[
g: \mathbf {B} (\mathbf {A} \otimes_ {k} \mathbf {A}, k) \rightarrow \mathbf {B} (\mathbf {A}, k) \otimes_ {k} \mathbf {B} (\mathbf {A}, k)
\]

such that

\[
g \left(\left[ x _ {1} \otimes y _ {1}, \dots , x _ {n} \otimes y _ {n} \right]\right) = \sum_ {0 \leqslant p \leqslant n} \pi \left(x _ {p + 1} \dots x _ {n}\right) \left[ x _ {1}, \dots , x _ {p} \right] \otimes y _ {1} \dots y _ {p} \left[ y _ {p + 1}, \dots , y _ {n} \right]
\]

for \(x_{1},\ldots,x_{n},y_{1},\ldots,y_{n}\) into A, is a morphism of resolutions.

\begin{enumerate}
\def\labelenumi{\alph{enumi})}
\setcounter{enumi}{1}
\tightlist
\item
  Assume A is commutative. Show that there exists on B(A, k) a structure of graded A-algebra such that one has for \(x_{1}, \ldots, x_{n} \in A\) :
\end{enumerate}

\[
\left[ x _ {1}, \dots , x _ {p} \right] \left[ x _ {p + 1}, \dots , x _ {n} \right] = \sum_ {\sigma \in M _ {p, n}} \varepsilon (\sigma) \left[ x _ {\sigma (1)}, \dots , x _ {\sigma (n)} \right]
\]

where \(M_{p,n}\) is the subset of \(S_{n}\) formed by the permutations \(\sigma\) such that \(\sigma(1) < \cdots < \sigma(p)\) and \(\sigma(p + 1) < \cdots < \sigma(n)\) .

Show that B(A, k) is then a differential graded A-algebra (X, p.~183, exercise 18).

\begin{enumerate}
\def\labelenumi{\arabic{enumi})}
\setcounter{enumi}{5}
\tightlist
\item
  Let P be a flat complex of right A-modules. For every A-module M and every integer i, we set
\end{enumerate}

\[
\mathrm{T} _ {i} ^ {\mathbf {P}} (\mathrm{M}) = \mathrm{H} _ {i} (\mathrm{P} \otimes_ {\mathbf {A}} \mathrm{M});
\]

if \(u:M\to N\) is an A-homomorphism, we denote by \(\mathbf{T}_{i}^{\mathbb{P}}(u)\) the \(k\) -homomorphism \(\mathbf{H}_{i}(1_{\mathbb{P}}\otimes u)\) .

\begin{enumerate}
\def\labelenumi{\alph{enumi})}
\tightlist
\item
  Let r be an integer. Show that the following conditions are equivalent:
\end{enumerate}

α) For every injection \(u: \mathbf{M}' \to \mathbf{M}\) , the map \(\mathbf{T}_{r}^{\mathbb{P}}(u)\) is injective.

β) For every surjection \(v: \mathbf{M} \to \mathbf{M}''\) , the map \(\mathbf{T}_{r+1}^{\mathbf{P}}(v)\) is surjective.

\(\gamma)\) The right A-module \(\mathbf{P}_{r}/\mathbf{B}_{r}(\mathbf{P})\) is flat.

δ) There exists a flat complex P' of right A-modules, whose differential \(d_{r+1}\) is zero, and for every A-module M and every integer i a k-isomorphism \(\varphi_{i}^{\mathrm{M}} : \mathrm{T}_{i}^{\mathrm{P}}(\mathrm{M}) \to \mathrm{T}_{i}^{\mathrm{P}'}(\mathrm{M})\) such that one has

\[
\mathrm{T} _ {i} ^ {\mathrm{P} ^ {\prime}} (u) \circ \varphi_ {i} ^ {\mathrm{M}} = \varphi_ {i} ^ {\mathrm{N}} \circ \mathrm{T} _ {i} ^ {\mathrm{P}} (u)
\]

for every A-homomorphism \(u:M\to N\)

(To prove that \(\gamma\) ) implies \(\delta\) ), define \(\mathbf{P}'\) by setting \(\mathbf{P}_{r+1}' = \mathbf{Z}_{r+1}(\mathbf{P})\) and \(\mathbf{P}_r' = \mathbf{P}_r / \mathbf{B}_r(\mathbf{P})\) .)

\begin{enumerate}
\def\labelenumi{\alph{enumi})}
\setcounter{enumi}{1}
\tightlist
\item
  We also assume that A is Noetherian and that the A-modules \(P_{i}\) are finitely generated. Show that the conditions in a) are equivalent to the following condition:
\end{enumerate}

e) There exists an A-module Q and, for every A-module M, a k-isomorphism \(\psi^{M}: T_{r}^{P}(M) \to \text{Hom}_{A}(Q, M)\) such that \(\psi^{N} \circ T_{r}^{P}(u) = \text{Hom}(1, u) \circ \psi^{M}\) for every A-homomorphism \(u: M \to N\) .

(To prove that δ) implies ε), reduce to the case where \(d_{r+1} = 0\) and take for Q the dual of \(Z_{r}(\mathbf{P})\) .) c) Show that the result of b) remains valid under the hypothesis that A is Noetherian, \(H_{i}(\mathbf{P})\) finitely generated for every i and \(H_{i}(\mathbf{P}) = 0\) for \(i < i_{0}\) (using prop. 10, p.~56).

\begin{enumerate}
\def\labelenumi{\arabic{enumi})}
\setcounter{enumi}{6}
\tightlist
\item
  Let C be a complex of right A-modules and M a left A-module. The torsion product of C and M is the graded k-module \(\mathcal{T}or^{\mathbf{A}}(\mathbf{C},\mathbf{M})=\mathrm{H}(\mathbf{C}\otimes_{\mathbf{A}}\mathbf{L}(\mathbf{M}))\) . If \(f:C\to C'\) is a morphism of complexes and \(g:M\to M'\) a homomorphism of A-modules, we write \(\mathcal{T}or^{\mathbf{A}}(f,g)=\mathrm{H}(f\otimes\mathbf{L}(g))\) . a) If \(f:C\to C'\) is a homologism, \(\mathcal{T}or^{\mathbf{A}}(f,1)\) is an isomorphism; if C is flat and bounded on the right, the graded k-module \(\mathcal{T}or^{\mathbf{A}}(\mathbf{C},\mathbf{M})\) is isomorphic to \(\mathrm{H}(\mathbf{C}\otimes_{\mathbf{A}}\mathbf{M})\) .
\end{enumerate}

\begin{enumerate}
\def\labelenumi{\alph{enumi})}
\setcounter{enumi}{1}
\tightlist
\item
  If \(0 \to C' \to C \to C'' \to 0\) is an exact sequence of complexes of right A-modules and if M is an A-module, show that there exists an exact sequence
\end{enumerate}

\[
\dots \rightarrow \mathcal {T} o r _ {n} ^ {\mathrm{A}} (\mathrm{C} ^ {\prime}, \mathrm{M}) \rightarrow \mathcal {T} o r _ {n} ^ {\mathrm{A}} (\mathrm{C}, \mathrm{M}) \rightarrow \mathcal {T} o r _ {n} ^ {\mathrm{A}} (\mathrm{C} ^ {\prime \prime}, \mathrm{M}) \rightarrow \mathcal {T} o r _ {n - 1} ^ {\mathrm{A}} (\mathrm{C} ^ {\prime}, \mathrm{M}) \rightarrow \dots .
\]

If \(0 \to M' \xrightarrow{i} M \xrightarrow{p} M'' \to 0\) is an exact sequence of A-modules, show that one has an exact sequence:

\[
\dots \rightarrow \mathcal {T} o r _ {n} ^ {\mathrm{A}} (\mathrm{C}, \mathrm{M} ^ {\prime}) \rightarrow \mathcal {T} o r _ {n} ^ {\mathrm{A}} (\mathrm{C}, \mathrm{M}) \rightarrow \mathcal {T} o r _ {n} ^ {\mathrm{A}} (\mathrm{C}, \mathrm{M} ^ {\prime \prime}) \rightarrow \mathcal {T} o r _ {n - 1} ^ {\mathrm{A}} (\mathrm{C}, \mathrm{M} ^ {\prime}) \rightarrow \dots .
\]

(Note that there exists a homotopism from L(M'') to the cone of L(i)).

\begin{enumerate}
\def\labelenumi{\alph{enumi})}
\setcounter{enumi}{2}
\tightlist
\item
  Show that there exists a spectral sequence 'E (X, pp.~175–177, exercises 14–17) converging to \(\mathcal{T}or^{\mathrm{A}}(\mathrm{C},\mathrm{M})\) , such that \(E_{pq}^{2} = \mathrm{Tor}_{p}^{\mathrm{A}}(\mathrm{H}_{q}(\mathrm{C}),\mathrm{M})\) If moreover the complex C is bounded on the right, show that there exists a spectral sequence E converging to \(\mathcal{T}or^{\mathrm{A}}(\mathrm{C},\mathrm{M})\) , such that \(E_{pq}^{2} = H_{p}(\mathrm{Tor}_{q}^{\mathrm{A}}(\mathrm{C},\mathrm{M}))\) (consider the double complex \(\mathrm{C}\otimes_{\mathrm{A}}\mathrm{L}(\mathrm{M})\) ).
\end{enumerate}

\begin{enumerate}
\def\labelenumi{\arabic{enumi})}
\setcounter{enumi}{7}
\tightlist
\item
  Let A, B be two k-algebras (associative and unital), M a right A-module, P a left B-module, N an (A, B)-bimodule.
\end{enumerate}

\begin{enumerate}
\def\labelenumi{\alph{enumi})}
\tightlist
\item
  Show that there exist two spectral sequences of \(k\) -modules \({}^{1}\mathrm{E}\) and \({}^{2}\mathrm{E}\) , converging to the same \(k\) -graded module, such that
\end{enumerate}

\[
{ } ^ { 1 } \mathrm{E} _ { p q } ^ { 2 } = \operatorname{Tor} _ { p } ^ { \mathrm{A} } ( \mathrm{M} , \operatorname{Tor} _ { q } ^ { \mathrm{B} } ( \mathrm{N} , \mathrm{P} ) ) ; \quad { } ^ { 2 } \mathrm{E} _ { p q } ^ { 2 } = \operatorname{Tor} _ { p } ^ { \mathrm{B} } ( \operatorname{Tor} _ { q } ^ { \mathrm{A} } ( \mathrm{M} , \mathrm{N} ) , \mathrm{P} ) .
\]

(Consider the double complex \(L(M) \otimes_{A} (N \otimes_{B} L(P))\) .)

\begin{enumerate}
\def\labelenumi{\alph{enumi})}
\setcounter{enumi}{1}
\item
  If N is a flat B-module, show that the spectral sequence \({}^{2}E\) converges to Tor \(^{A}\) (M, N ⊗ \(_{B}\) P). If N is a flat A-module, the sequence \({}^{1}E\) converges to Tor \(^{B}\) (M ⊗ \(_{A}\) N, P).
\item
  Let \(\rho: \mathbf{A} \to \mathbf{B}\) a homomorphism of \(k\) -algebras. Deduce from \(b\) ) a spectral sequence E such that
\end{enumerate}

\[
\mathrm{E} _ {p q} ^ {2} = \operatorname{Tor} _ {p} ^ {\mathbf {B}} \left(\operatorname{Tor} _ {q} ^ {\mathbf {A}} (\mathbf {M}, \mathbf {B}), \mathbf {P}\right),
\]

converging to Tor \(^{A}\) (M, P).

\begin{enumerate}
\def\labelenumi{\arabic{enumi})}
\setcounter{enumi}{8}
\tightlist
\item
  Let P be a complex of right A-modules, Q a complex of left A-modules. Suppose either that P and Q are bounded below, or that there exists an integer d such that every left or right A-module admits a projective resolution of length \(\leqslant d\) .
\end{enumerate}

\begin{enumerate}
\def\labelenumi{\alph{enumi})}
\tightlist
\item
  Show that there exist two spectral sequences (X, p.~175, exercise 14) 'E and ``E, converging to the same graded k-module, such that
\end{enumerate}

\[
^ {\prime} \mathrm{E} _ {p q} ^ {2} = \mathrm{H} _ {p} (\operatorname{Tor} _ {q} ^ {\mathbf {A}} (\mathrm{P}, \mathrm{Q})); \quad^ {\prime \prime} \mathrm{E} _ {p q} ^ {2} = \bigoplus_ {q ^ {\prime} + q ^ {\prime \prime} = q} \operatorname{Tor} _ {p} ^ {\mathbf {A}} \left(\mathrm{H} _ {q ^ {\prime}} (\mathrm{P}), \mathrm{H} _ {q ^ {\prime \prime}} (\mathrm{Q})\right)
\]

\((\mathrm{where}\ \mathrm{Tor}_{q}^{\mathrm{A}}(\mathrm{P},\mathrm{Q})\ \mathrm{denotes}\ \mathrm{the}\ \mathrm{complex}\ \mathrm{associated}\ \mathrm{with}\ \mathrm{the}\ \mathrm{bicomplexe}\ \oplus\ \mathrm{Tor}_{q}^{\mathrm{A}}(\mathrm{P}_{r},\mathrm{Q}_{s}))\)

(Consider projective resolutions (X, p.~183, exercise 17) \(\mathcal{P}\) and \(\mathcal{Q}\) of P and Q, and the double complex C such that \(C_{p,q} = \bigoplus (\mathcal{P}_{p'q'} \otimes \mathcal{Q}_{p''q''})\) .)

\begin{enumerate}
\def\labelenumi{\alph{enumi})}
\setcounter{enumi}{1}
\tightlist
\item
  If P is flat, show that the spectral sequence ``E converges to H(P ⊗\_A Q). If H(P) is flat, the sequence 'E converges to H(P) ⊗\_A H(Q).
\end{enumerate}

\exercisesection{Extension Modules}

¶ 1) An A-complex C is said to be pure if for every right A-module P one has H(P ⊗\_A C) = 0.

\begin{enumerate}
\def\labelenumi{\alph{enumi})}
\item
  A complex homotopic to zero is pure; a complex of flat modules, with zero homology and bounded on the right, is pure. Give an example of a complex of free modules of rank 1, with zero homology but not pure.
\item
  Show that the following conditions are equivalent:
\end{enumerate}

α) C is pure.

β) For every finitely presented A-module M, one has H(Homgr\_A (M, C)) = 0.

\(\gamma)\) There exists a filtered inductive system \((C^{(\alpha)})_{\alpha \in j}\) of complexes homotopic to zero such that \(C = \varinjlim C^{(\alpha)}\) .

(To see that \(\beta\) ) implies \(\gamma\) ), reduce to the case where \(C_i = 0\) for \(i < i_0\) and write \(C_{i_0}\) as an inductive limit of finitely presented modules.)

\begin{enumerate}
\def\labelenumi{\arabic{enumi})}
\setcounter{enumi}{1}
\tightlist
\item
  Let M be an A-module. We say that a submodule \(\mathbf{M}'\) of M is pure if the complex defined by the exact sequence \(0 \to \mathbf{M}' \to \mathbf{M} \to \mathbf{M} / \mathbf{M}' \to 0\) is pure (Exercise 1).
\end{enumerate}

\begin{enumerate}
\def\labelenumi{\alph{enumi})}
\item
  Show that when A is a principal ideal ring, the notion of pure submodule coincides with that of VII, § 2, Exercise 7.
\item
  Suppose that the A-module M is flat. Show that the following conditions are equivalent: \(\alpha\) the submodule M' is pure;
\end{enumerate}

β) the quotient M/M' is flat;

\(\gamma)\) for every ideal \(\mathfrak{a}\) of \(\mathbf{A}\) , we have \(\mathbf{M}' \cap \mathfrak{a}\mathbf{M} = \mathfrak{a}\mathbf{M}'\) .

\begin{enumerate}
\def\labelenumi{\arabic{enumi})}
\setcounter{enumi}{2}
\tightlist
\item
  Let M be an A-module, n an integer ≥ 1. Show that the following conditions are equivalent:
\end{enumerate}

α) M admits a finite n-presentation (X, p.~180, Exercise 6).

β) For every filtered inductive system \((\mathrm{N}_{\alpha})_{\alpha\in I}\) , the canonical homomorphism

\[
\varinjlim \operatorname{Ext} _ {\mathbf {A}} ^ {i} (\mathbf {M}, \mathbf {N} _ {\alpha}) \rightarrow \operatorname{Ext} _ {\mathbf {A}} ^ {i} (\mathbf {M}, \varinjlim \mathbf {N} _ {\alpha})
\]

is bijective for i \textless{} n.

γ) For every family \((\mathbf{P}_{\beta})_{\beta \in J}\) of right A-modules, the canonical homomorphism

\[
\operatorname{Tor} _ {i} ^ {\mathbf {A}} \left(\prod_ {\beta \in J} \mathrm{P} _ {\beta}, \mathrm{M}\right)\rightarrow \prod_ {\beta \in J} \operatorname{Tor} _ {i} ^ {\mathbf {A}} \left(\mathrm{P} _ {\beta}, \mathrm{M}\right) \quad \text { is   bijective   for } \quad i <   n.
\]

δ) For every set I, the canonical homomorphism \(A_{d}^{I} \otimes_{A} M \to M^{I}\) is bijective, and we have

\[
\operatorname{Tor} _ {i} ^ {\mathbf {A}} \left(\mathrm{A} _ {d} ^ {\mathrm{I}}, \mathrm{M}\right) = 0 \quad \text { for } \quad 0 <   i <   n.
\]

(Argue by induction on n, using Exercise 18, p.~170.)

\begin{enumerate}
\def\labelenumi{\arabic{enumi})}
\setcounter{enumi}{3}
\tightlist
\item
  Let P be a complex of projective A-modules, zero on the right, and n an integer such that \(H_{i}(P)=0\) for i\textless n.~Let M be a left A-module, N a right A-module. Show that we have \(H^{i}(\mathrm{Homgr}_{\mathbf{A}}(\mathbf{P},\mathbf{M}))=0\) and \(H_{i}(\mathbf{N}\otimes_{\mathbf{A}}\mathbf{P})=0\) for i\textless n, and that the canonical homomorphisms
\end{enumerate}

\[
\lambda^ {n}: \mathrm{H} ^ {n} (\operatorname{Homgr} _ {\mathbf {A}} (\mathrm{P}, \mathrm{M})) \rightarrow \operatorname{Hom} _ {\mathbf {A}} \left(\mathrm{H} _ {n} (\mathrm{P}), \mathrm{M}\right), \quad \gamma^ {n}: \mathrm{N} \otimes_ {\mathbf {A}} \mathrm{H} _ {n} (\mathrm{P}) \rightarrow \mathrm{H} _ {n} (\mathrm{N} \otimes_ {\mathbf {A}} \mathrm{P})
\]

are bijective.

¶ 5) Let P be a projective A-complex, zero on the right, n an integer ≥ 0. Show that the following conditions are equivalent:

α) There exists a complex \(\mathbf{P}'\) of finitely generated projective A-modules, such that \(\mathbf{P}_i' = 0\) for \(i > n\) or \(i < 0\) , and a morphism \(f: \mathbf{P}' \to \mathbf{P}\) such that \(\mathbf{H}_i(f)\) is bijective for \(i < n\) and surjective for \(i = n\) .

β) For every filtered inductive system of A-modules \((M_{\alpha})_{\alpha \in I}\) , the canonical homomorphism

\[
\varinjlim \mathrm{H} ^ {i} (\operatorname{Homgr} _ {\mathrm{A}} (\mathrm{P}, \mathrm{M} _ {\alpha})) \rightarrow \mathrm{H} ^ {i} (\operatorname{Homgr} _ {\mathrm{A}} (\mathrm{P}, \varinjlim \mathrm{M} _ {\alpha}))
\]

is bijective for \(i < n\) and injective for \(i = n\) .

γ) For any family of right A-modules \((\mathbf{N}_{\alpha})_{\alpha\in I}\) , the canonical homomorphism

\[
\mathrm{H} _ {i} ((\prod_ {\alpha} \mathrm{N} _ {\alpha}) \otimes_ {\mathrm{A}} \mathrm{P}) \rightarrow \prod_ {\alpha} \mathrm{H} _ {i} (\mathrm{N} _ {\alpha} \otimes \mathrm{P})
\]

is bijective for \(i < n\) and surjective for \(i = n\) .

(To prove that \(\alpha\) ) implies \(\beta\) ) and \(\gamma\) ), apply Exercise 4 to the cone of \(f\). To prove that \(\beta\) )

or γ) implies α), reason by induction on n, applying Exercise 4, Exercise 18, p.~170 and Exercise 10, p.~174.)

\begin{enumerate}
\def\labelenumi{\arabic{enumi})}
\setcounter{enumi}{5}
\tightlist
\item
  Let A, B be two k-algebras (associative and unital), M a left A-module, P a left B-module, N a (B, A)-bimodule.
\end{enumerate}

\begin{enumerate}
\def\labelenumi{\alph{enumi})}
\tightlist
\item
  Show that there exist two spectral sequences (X, pp.~175-177, exercises 14 to 17) of \(k\) -modules \({}^{1}\mathrm{E}\) and \({}^{2}\mathrm{E}\) , converging to the same \(k\) -graded module, such that
\end{enumerate}

\[
{ } ^ { 1 } \mathrm{E} _ { 2 } ^ { p q } = \operatorname{Ext} _ { \mathbf { A } } ^ { p } ( \mathbf { M } , \operatorname{Ext} _ { \mathbf { B } } ^ { q } ( \mathbf { N } , \mathbf { P } ) ) ; \quad { } ^ { 2 } \mathrm{E} _ { 2 } ^ { p q } = \operatorname{Ext} _ { \mathbf { B } } ^ { p } ( \operatorname{Tor} _ { q } ^ { \mathbf { A } } ( \mathbf { N } , \mathbf { M } ) , \mathbf { P } ) .
\]

\begin{enumerate}
\def\labelenumi{\alph{enumi})}
\setcounter{enumi}{1}
\tightlist
\item
  By composing the homomorphisms \(\gamma\) and \(\delta\) of Exercise 15, p.~176, define a \(k\) -homomorphism
\end{enumerate}

\[
\mu : \operatorname{Ext} _ {\mathbf {A}} (\mathbf {M}, \operatorname{Hom} _ {\mathbf {B}} (\mathbf {N}, \mathbf {P})) \rightarrow \operatorname{Hom} _ {\mathbf {B}} \left(\operatorname{Tor} ^ {\mathbf {A}} (\mathbf {N}, \mathbf {M}), \mathbf {P}\right).
\]

If P is injective, prove that μ is an isomorphism.

\begin{enumerate}
\def\labelenumi{\alph{enumi})}
\setcounter{enumi}{2}
\item
  If N is a flat A-module, show that the sequence \({}^{1}E\) converges to Ext \(_{B}(N \otimes_{A} M, P)\) ; if N is a projective B-module, the sequence \({}^{2}E\) converges to Ext \(_{A}(M, Hom_{B}(N, P))\) .
\item
  Let \(\rho: A \to B\) a homomorphism of \(k\) -algebras. Show that there exists a spectral sequence E, such that \(\mathbf{E}_2^{pq} = \operatorname{Ext}_{\mathbf{B}}^p (\operatorname{Tor}_q^A(B, M), P)\) , converging to \(\operatorname{Ext}_{\mathbf{A}}(M, P)\) and a spectral sequence 'E, such that
\end{enumerate}

\[
{ } ^ { \prime } \mathrm{E} _ { 2 } ^ { p q } = \operatorname{Ext} _ { \mathbf { B } } ^ { p } ( \mathbf { P } , \operatorname{Ext} _ { \mathbf { A } } ^ { q } ( \mathbf { B } , \mathbf { M } ) ) ,
\]

converging to \(\operatorname{Ext}_{\mathbf{A}}(\mathbf{P}, \mathbf{M})\) .

\begin{enumerate}
\def\labelenumi{\arabic{enumi})}
\setcounter{enumi}{6}
\tightlist
\item
  Let C be an A-complex and M an A-module. One calls the module of extensions of M by C (resp. of C by M) the graded k-module
\end{enumerate}

\[
\mathcal {E} x t _ {\mathrm{A}} (\mathrm{C}, \mathrm{M}) = \mathrm{H} (\operatorname{Homgr} _ {\mathrm{A}} (\mathrm{C}, \mathrm{I} (\mathrm{M}))) \quad \left(\text { resp. } \mathcal {E} x t _ {\mathrm{A}} (\mathrm{M}, \mathrm{C}) = \mathrm{H} (\operatorname{Homgr} _ {\mathrm{A}} (\mathrm{L} (\mathrm{M}), \mathrm{C}))\right).
\]

If \(f: \mathbf{C} \to \mathbf{C}'\) is a morphism of complexes and \(g: \mathbf{M} \to \mathbf{M}'\) a homomorphism of A-modules, we set \(\mathcal{E}xt_{\mathbf{A}}(f, g) = \mathrm{H}(\mathrm{Homgr}(f, \mathrm{l}(g)))\) , \(\mathcal{E}xt_{\mathbf{A}}(g, f) = \mathrm{H}(\mathrm{Homgr}(\mathrm{L}(g), f))\) .

\begin{enumerate}
\def\labelenumi{\alph{enumi})}
\item
  If \(f: \mathbf{C} \to \mathbf{C}'\) is a homologism, the \(k\) -homomorphisms \(\mathcal{E}xt_{\mathbf{A}}(f, 1)\) and \(\mathcal{E}xt_{\mathbf{A}}(1, f)\) are bijective. If C is projective and bounded on the right (resp. injective and bounded on the left), the \(k\) -graded module \(\mathcal{E}xt_{\mathbf{A}}(\mathbf{C}, \mathbf{M})\) (resp. \(\mathcal{E}xt_{\mathbf{A}}(\mathbf{M}, \mathbf{C})\) ) is isomorphic to \(\mathrm{H}(\mathrm{Homgr}_{\mathbf{A}}(\mathbf{C}, \mathbf{M}))\) (resp. \(\mathrm{H}(\mathrm{Homgr}_{\mathbf{A}}(\mathbf{M}, \mathbf{C}))\)) .
\item
  Let \(0 \to C' \to C \to C'' \to 0\) be an exact sequence of A-complexes and \(0 \to M' \to M \to M'' \to 0\) an exact sequence of A-modules. Show that there exist exact sequences
\end{enumerate}

\[
\begin{array}{l} \dots \to \mathcal {E} x t _ {\mathrm{A}} ^ {n} (\mathrm{C} ^ {\prime \prime}, \mathrm{M}) \to \mathcal {E} x t _ {\mathrm{A}} ^ {n} (\mathrm{C}, \mathrm{M}) \to \mathcal {E} x t _ {\mathrm{A}} ^ {n} (\mathrm{C} ^ {\prime}, \mathrm{M}) \to \mathcal {E} x t _ {\mathrm{A}} ^ {n + 1} (\mathrm{C} ^ {\prime \prime}, \mathrm{M}) \to \dots \\ \dots \to \mathcal {E} x t _ {\mathrm{A}} ^ {n} (\mathrm{C}, \mathrm{M} ^ {\prime}) \to \mathcal {E} x t _ {\mathrm{A}} ^ {n} (\mathrm{C}, \mathrm{M}) \to \mathcal {E} x t _ {\mathrm{A}} ^ {n} (\mathrm{C}, \mathrm{M} ^ {\prime \prime}) \to \mathcal {E} x t _ {\mathrm{A}} ^ {n + 1} (\mathrm{C}, \mathrm{M} ^ {\prime}) \to \dots \\ \dots \to \mathcal {E} x t _ {\mathrm{A}} ^ {n} (\mathrm{M}, \mathrm{C} ^ {\prime}) \to \mathcal {E} x t _ {\mathrm{A}} ^ {n} (\mathrm{M}, \mathrm{C}) \to \mathcal {E} x t _ {\mathrm{A}} ^ {n} (\mathrm{M}, \mathrm{C} ^ {\prime \prime}) \to \mathcal {E} x t _ {\mathrm{A}} ^ {n + 1} (\mathrm{M}, \mathrm{C} ^ {\prime}) \to \dots \\ \dots \to \mathcal {E} x t _ {\mathrm{A}} ^ {n} (\mathrm{M} ^ {\prime \prime}, \mathrm{C}) \to \mathcal {E} x t _ {\mathrm{A}} ^ {n} (\mathrm{M}, \mathrm{C}) \to \mathcal {E} x t _ {\mathrm{A}} ^ {n} (\mathrm{M} ^ {\prime}, \mathrm{C}) \to \mathcal {E} x t _ {\mathrm{A}} ^ {n + 1} (\mathrm{M} ^ {\prime \prime}, \mathrm{C}) \to \dots \end{array}
\]

\begin{enumerate}
\def\labelenumi{\alph{enumi})}
\setcounter{enumi}{2}
\tightlist
\item
  Show that there exists a spectral sequence ``E (X, pp.~175-177, exercises 14 to 17) converging to \(\mathcal{E}xt_{\mathrm{A}}(\mathrm{C},\mathrm{M})\) such that \(E_2^{pq} = H^q (\mathrm{Ext}_A^p (\mathrm{C},\mathrm{M}))\) . If C is bounded above, show that there exists a spectral sequence 'E, converging to \(\mathcal{E}xt_{\mathrm{A}}(\mathrm{C},\mathrm{M})\) , such that \(E_2^{pq} = \mathrm{Ext}_A^p (\mathrm{H}_q(\mathrm{C}),\mathrm{M})\) (consider the bicomplex \(\mathrm{Homgr}_{\mathrm{A}}(\mathrm{C},\mathrm{I}(\mathrm{M}))\) ).
\end{enumerate}

\begin{enumerate}
\def\labelenumi{\arabic{enumi})}
\setcounter{enumi}{7}
\tightlist
\item
  Let P, Q be two A-complexes. Suppose either that P is bounded below and Q bounded above, or that there exists an integer d such that every A-module admits a projective resolution of length \(\leqslant d\) .
\end{enumerate}

\begin{enumerate}
\def\labelenumi{\alph{enumi})}
\tightlist
\item
  Show that there exist two spectral sequences 'E and ``E, converging to the same graded k-module, such that
\end{enumerate}

\[
^ {\prime} \mathrm{E} _ {2} ^ {p q} = \mathrm{H} ^ {p} (\mathrm{Ext} _ {\mathrm{A}} ^ {q} (\mathrm{P}, \mathrm{Q})) \quad^ {\prime \prime} \mathrm{E} _ {2} ^ {p q} = \bigoplus_ {q ^ {\prime} + q ^ {\prime \prime} = q} \mathrm{Ext} _ {\mathrm{A}} ^ {p} (\mathrm{H} _ {q ^ {\prime}} (\mathrm{P}), \mathrm{H} ^ {q ^ {\prime \prime}} (\mathrm{Q}))
\]

\((\text{ou} \operatorname{Ext}_{\mathbf{A}}^{q}(\mathbf{P}, \mathbf{Q}) \text{ denotes the complex associated with the bicomplex } \oplus \operatorname{Ext}_{\mathbf{A}}^{q}(\mathbf{P}_{r}, \mathbf{Q}^{s}))\) .

\begin{enumerate}
\def\labelenumi{\alph{enumi})}
\setcounter{enumi}{1}
\tightlist
\item
  If P is projective or Q is injective, the spectral sequence ``E converges to H(Homgr \(_A\) (P, Q)). If H(P) is projective or H(Q) is injective, the spectral sequence 'E converges to Homgr \(_A\) (H(P), H(Q)).
\end{enumerate}

\exercisesection{Use of Non-Canonical Resolutions}

\begin{enumerate}
\def\labelenumi{\arabic{enumi})}
\tightlist
\item
  Let \(\rho: A \to B\) be a ring homomorphism, M, N two left B-modules, P a right B-module.
\end{enumerate}

\begin{enumerate}
\def\labelenumi{\alph{enumi})}
\tightlist
\item
  Let \(L_{B}\) (resp. \(L_{A}\) ) be a projective resolution of the B-module (resp. of the A-module) M. Show that there exists an A-linear morphism of resolutions \(L_{A} \rightarrow L_{B}\) , unique up to homotopy; deduce graded k-homomorphisms of degree zero
\end{enumerate}

\[
\alpha : \operatorname{Tor} ^ {\mathrm{A}} (\mathrm{P}, \mathrm{M}) \rightarrow \operatorname{Tor} ^ {\mathrm{B}} (\mathrm{P}, \mathrm{M}) \quad \beta : \operatorname{Ext} _ {\mathrm{B}} (\mathrm{M}, \mathrm{N}) \rightarrow \operatorname{Ext} _ {\mathrm{A}} (\mathrm{M}, \mathrm{N}).
\]

\begin{enumerate}
\def\labelenumi{\alph{enumi})}
\setcounter{enumi}{1}
\tightlist
\item
  Consider the canonical homomorphisms (X, p.~109)
\end{enumerate}

\[
\begin{array}{l l} \lambda : \operatorname{Tor} ^ {\mathsf {A}} (\mathsf {P}, \mathsf {M}) \to \operatorname{Tor} ^ {\mathsf {B}} (\mathsf {P} \otimes_ {\mathsf {A}} \mathsf {B}, \mathsf {M}) & \mu : \operatorname{Tor} ^ {\mathsf {A}} (\mathsf {P}, \mathsf {M}) \to \operatorname{Tor} ^ {\mathsf {B}} (\mathsf {P}, \mathsf {B} \otimes_ {\mathsf {A}} \mathsf {M}) \\ v: \operatorname{Ext} _ {\mathsf {B}} (\mathsf {B} \otimes_ {\mathsf {A}} \mathsf {M}, \mathsf {N}) \to \operatorname{Ext} _ {\mathsf {A}} (\mathsf {M}, \mathsf {N}) & \rho : \operatorname{Ext} _ {\mathsf {B}} (\mathsf {M}, \operatorname{Hom} _ {\mathsf {A}} (\mathsf {B}, \mathsf {N})) \to \operatorname{Ext} _ {\mathsf {A}} (\mathsf {M}, \mathsf {N}). \end{array}
\]

Let \(m: \mathbf{B} \otimes_{\mathbf{A}} \mathbf{M} \to \mathbf{M}\) , \(p: \mathbf{P} \otimes_{\mathbf{A}} \mathbf{B} \to \mathbf{B}\) and \(n: \mathbf{N} \to \operatorname{Hom}_{\mathbf{A}}(\mathbf{B}, \mathbf{N})\) the B-homomorphisms defined by \(m(b \otimes x) = bx\) , \(p(y \otimes b) = yb\) and \(n(z)(b) = bz\) for \(b \in \mathbf{B}\) , \(x \in \mathbf{M}\) , \(y \in \mathbf{P}\) , \(z \in \mathbf{N}\) . Prove the equalities

\[
\begin{array}{l} \alpha = \operatorname{Tor} (1 _ {\mathrm{P}}, m) \circ \mu = \operatorname{Tor} (p, 1 _ {\mathrm{M}}) \circ \lambda \\ \beta = v \circ \operatorname{Ext} (m, 1 _ {\mathrm{N}}) = \rho \circ \operatorname{Ext} (1 _ {\mathrm{M}}, n). \end{array}
\]

\begin{enumerate}
\def\labelenumi{\arabic{enumi})}
\setcounter{enumi}{1}
\tightlist
\item
  Let A, B be two k-algebras (associative and unital); let M be a left A-module, P a right B-module, N an (A, B)-bimodule. Consider the homomorphism
\end{enumerate}

\[
\sigma : \operatorname{Hom} _ {\mathrm{B}} (\mathrm{N}, \mathrm{P}) \otimes_ {\mathrm{A}} \mathrm{M} \rightarrow \operatorname{Hom} _ {\mathrm{B}} (\operatorname{Hom} _ {\mathrm{A}} (\mathrm{M}, \mathrm{N}), \mathrm{P})
\]

such that \(\sigma(u \otimes m)(v) = u \circ v(m)\) for \(u \in \mathrm{Hom}_{\mathbf{B}}(\mathbf{N}, \mathbf{P})\) , \(v \in \mathrm{Hom}_{\mathbf{A}}(\mathbf{M}, \mathbf{N})\) , \(m \in \mathbf{M}\) (II, p.~190, exercise 6). a) By replacing M by a projective resolution, define using \(\sigma\) a \(k\) -homomorphism

\[
\tau : \operatorname{Tor} ^ {\mathbf {A}} \left(\operatorname{Hom} _ {\mathbf {B}} (\mathrm{N}, \mathrm{P}), \mathrm{M}\right)\rightarrow \operatorname{Homgr} _ {\mathbf {B}} \left(\operatorname{Ext} _ {\mathbf {A}} (\mathrm{M}, \mathrm{N}), \mathrm{P}\right).
\]

\begin{enumerate}
\def\labelenumi{\alph{enumi})}
\setcounter{enumi}{1}
\item
  Assume henceforth that A is left noetherian and that the A-module M is finitely generated. Show that if the B-module P is injective, \(\tau\) is an isomorphism.
\item
  Show that there exist two spectral sequences E and 'E (X, pp.~175-177, exercises 14 to 17), converging to the same graded k-module, such that
\end{enumerate}

\[
E _ {2} ^ {p q} = \operatorname{Tor} _ {- p} ^ {A} \left(\operatorname{Ext} _ {B} ^ {q} (N, P), M\right) \quad^ {\prime} E _ {2} ^ {p q} = \operatorname{Ext} _ {B} ^ {p} \left(\operatorname{Ext} _ {A} ^ {- q} (M, N), P\right).
\]

\begin{enumerate}
\def\labelenumi{\arabic{enumi})}
\setcounter{enumi}{2}
\item
  Let I be an injective \(k\) -module; for every A-module M, denote by D(M) the right A-module \(\mathrm{Hom}_{k}(M, I)\) . Let M and N be two left A-modules, P a right A-module. Prove that for every \(p \geqslant 0\) , the \(k\) -modules \(\mathrm{Ext}_{A}^{p}(P, D(M))\) and D( \(\mathrm{Tor}_{p}^{A}(P, M)\) ) are isomorphic; if A is left noetherian and if M is a finitely generated A-module, the \(k\) -modules \(\mathrm{Tor}_{p}^{A}(D(N), M)\) and D( \(\mathrm{Ext}_{A}^{p}(M, N)\) ) are isomorphic (use exercise 2, b) and exercise 6, b), p.~188).
\item
  For every group G, denote by \(\varepsilon_{\mathbf{G}}: \mathbf{Z}^{(\mathbf{G})} \to \mathbf{Z}\) the ring homomorphism such that \(\varepsilon_{\mathbf{G}}(e_g) = 1\) for all \(g \in \mathbf{G}\) ; set \(\mathrm{I}_{\mathbf{G}} = \operatorname{Ker}(\varepsilon_{\mathbf{G}})\) .
\end{enumerate}

\begin{enumerate}
\def\labelenumi{\alph{enumi})}
\item
  Show that the group \(\mathbf{H}_{1}(\mathbf{G},\mathbf{Z})\) is isomorphic to \(\mathrm{I}_{\mathrm{G}} / \mathrm{I}_{\mathrm{G}}^{2}\) .
\item
  Let \((g_{\iota})_{\iota \in \mathrm{I}}\) be a generating family of G; show that the elements \((e_{g_{\iota}} - 1)\) , for \(\iota \in \mathrm{I}\) , generate the ideal \(\mathbf{I}_{\mathbf{G}}\) .
\item
  Suppose that the family \((g_{t})_{t\in I}\) is basic (I, p.~86). Show that \((e_{g_t} - 1)_{t\in I}\) is a basis of the left \(\mathbf{Z}^{(G)}\) -module \(I_G\) .
\end{enumerate}

\begin{enumerate}
\def\labelenumi{\arabic{enumi})}
\setcounter{enumi}{4}
\item
  Let X be a set, \(G = F(X)\) the free group constructed on X, M a \(Z^{(G)}\) -module. Show that the groups \(H^q(G, M)\) and \(H_q(G, M)\) are zero for \(q \geqslant 2\) (use exercise 4).
\item
  Let G be a finite group of order n, M a \(\mathbf{Z}^{(G)}\) -module. Show that for every \(q \geqslant 1\) , the Z-modules \(\mathrm{H}^{q}(\mathrm{G}, \mathrm{M})\) and \(\mathrm{H}_{q}(\mathrm{G}, \mathrm{M})\) are annihilated by n (consider the homotopy h of C(G, M) into itself defined by
\end{enumerate}

\[
h f (g _ {1}, \dots , g _ {q}) = (- 1) ^ {q} \sum_ {g \in G} f (g _ {1}, \dots , g _ {q}, g)
\]

for \(f \in \mathbf{C}^{q+1}(\mathbf{G}, \mathbf{M})\) and \(g, g_1, \ldots, g_q \in \mathbf{G}\) ). If \(\mathbf{M}\) is a flat \(\mathbf{Z}\) finitely generated module, the groups \(\mathrm{H}^q(\mathbf{G}, \mathbf{M})\) and \(\mathrm{H}_q(\mathbf{G}, \mathbf{M})\) are finite for \(q \geqslant 1\) .

\begin{enumerate}
\def\labelenumi{\arabic{enumi})}
\setcounter{enumi}{6}
\tightlist
\item
  Let G be a cyclic group of order n, M a \(\mathbf{Z}^{(G)}\) -module, \(\sigma\) a generator of G. Denote by N the multiplication in M by the element \(\sum_{g\in G}e_{g}\) of \(\mathbf{Z}^{(G)}\) . Show that the Z-module \(\mathrm{H}^{0}(\mathrm{G},\mathrm{M})\) (resp. \(\mathrm{H}_{0}(\mathrm{G},\mathrm{M})\) ) is iso-
\end{enumerate}

morphic to Ker \((1 - \sigma_{\mathbf{M}})\) (resp. Coker \((1 - \sigma_{\mathbf{M}})\) ); for \(q\) odd, the groups \(\mathrm{H}^q (\mathrm{G},\mathrm{M})\) and \(\mathrm{H}_{q + 1}(\mathrm{G},\mathrm{M})\) are isomorphic to Ker N/Im \((1 - \sigma_{\mathbf{M}})\) ; for \(q\) even \(\geqslant 2\) , \(\mathrm{H}^q (\mathrm{G},\mathrm{M})\) and \(\mathrm{H}_{q - 1}(\mathrm{G},\mathrm{M})\) are isomorphic to Ker \((1 - \sigma_{\mathbf{M}}) / \mathrm{Im}\) N. (Use exercise 1, p.~178.)

\begin{enumerate}
\def\labelenumi{\arabic{enumi})}
\setcounter{enumi}{7}
\item
  Let F be a free group, R a normal subgroup of F, and G = F/R. Show that the group H₂(G, Z) is isomorphic to ((F, F) ∩ R)/(F, R) (use Exercise 3, p.~179).
\item
  \begin{enumerate}
  \def\labelenumii{\alph{enumii})}
  \tightlist
  \item
    Let G be a group, M a G-module, and g an element of the center of G such that the endomorphism \(m \mapsto gm - m\) of M is bijective. Show that the groups \(\mathrm{H}^{p}(\mathrm{G}, \mathrm{M})\) and \(\mathrm{H}_{p}(\mathrm{G}, \mathrm{M})\) are zero for \(p \geqslant 1\) . b) Let V be a vector space over a field k, with Card(k) \(\geqslant 3\) . Prove that the groups \(\mathrm{H}^{p}(\mathrm{GL}(\mathrm{V}), \mathrm{V})\) and \(\mathrm{H}_{p}(\mathrm{GL}(\mathrm{V}), \mathrm{V})\) are zero for \(p \geqslant 1\) .
  \end{enumerate}
\item
  Let G and H be two groups, \(u: H \to G\) a homomorphism, M a \(\mathbf{Z}^{(G)}\) -module. For every \(q \geqslant 0\) , we deduce from the homomorphisms \(\alpha\) and \(\beta\) of Exercise 1 homomorphisms
\end{enumerate}

\[
u _ {q} ^ {\mathrm{M}}: \mathrm{H} _ {q} (\mathrm{H}, \mathrm{M}) \rightarrow \mathrm{H} _ {q} (\mathrm{G}, \mathrm{M}) \quad u _ {\mathrm{M}} ^ {q}: \mathrm{H} ^ {q} (\mathrm{G}, \mathrm{M}) \rightarrow \mathrm{H} ^ {q} (\mathrm{H}, \mathrm{M})
\]

which we simply denote by \(u_q\) and \(u^q\) if there is no ambiguity about M.

\begin{enumerate}
\def\labelenumi{\alph{enumi})}
\tightlist
\item
  For \(f \in \mathbf{C}^q(\mathbf{G}, \mathbf{M})\) , we denote by \(f_u\) the element of \(\mathbf{C}^q(\mathbf{H}, \mathbf{M})\) such that \(f_u(h_1, \ldots, h_q) = f(u(h_1), \ldots, u(h_q))\) for \(h_1, \ldots, h_q \in \mathbf{H}\) . Show that the map \(f \mapsto f_u\) is a morphism of complexes from \(\mathbf{C}'(\mathbf{G}, \mathbf{M})\) to \(\mathbf{C}'(\mathbf{H}, \mathbf{M})\) , which induces on homology the homomorphism \(\bigoplus_{q} u^q\) . Define similarly a morphism from \(\mathbf{C}'(\mathbf{H}, \mathbf{M})\)
\end{enumerate}

into C'(G, M) which induces ( \(\oplus u_{q}\) ).

\begin{enumerate}
\def\labelenumi{\alph{enumi})}
\setcounter{enumi}{1}
\tightlist
\item
  Assume H = G and take for u an inner automorphism of G. Show that the homomorphisms \(u_{q}\) and \(u^{q}\) are equal to the identity (argue by induction on q).
\end{enumerate}

\begin{enumerate}
\def\labelenumi{\arabic{enumi})}
\setcounter{enumi}{10}
\tightlist
\item
  Let G be a group, H a distinguished subgroup of G, M a G-module. By Exercise 10, the group G acts by conjugation on the Z-modules H \(^{q}\) (H, M); since the action of the subgroup H is trivial, these are therefore endowed with a structure of Z \(^{(G/H)}\) -modules.
\end{enumerate}

\begin{enumerate}
\def\labelenumi{\alph{enumi})}
\tightlist
\item
  Show that there exists a spectral sequence E (X, p.~175, Exercise 14), converging to \(\bigoplus_{n} H^{n}(G, M)\) ,
\end{enumerate}

such that \(E_{2}^{pq} = H^{p}(G/H, H^{q}(H, M))\) (consider the spectral sequence of X, p.~188, Exercise 6, d), taking into account the canonical isomorphism (12), p.~109).

\begin{enumerate}
\def\labelenumi{\alph{enumi})}
\setcounter{enumi}{1}
\tightlist
\item
  Show that there exists an exact sequence
\end{enumerate}

\[
0 \to \mathrm{H} ^ {1} (\mathrm{G} / \mathrm{H}, \mathrm{H} ^ {0} (\mathrm{H}, \mathrm{M})) \to \mathrm{H} ^ {1} (\mathrm{G}, \mathrm{M}) \to \mathrm{H} ^ {0} (\mathrm{G} / \mathrm{H}, \mathrm{H} ^ {1} (\mathrm{H}, \mathrm{M})) \to \mathrm{H} ^ {2} (\mathrm{G} / \mathrm{H}, \mathrm{H} ^ {0} (\mathrm{H}, \mathrm{M})) \to \mathrm{H} ^ {2} (\mathrm{G}, \mathrm{M})
\]

Describe the homomorphisms of this sequence using the description of H'(G, M) in terms of C'(G, M). c) Define a spectral sequence analogous to that of a) for homology.

\begin{enumerate}
\def\labelenumi{\arabic{enumi})}
\setcounter{enumi}{11}
\tightlist
\item
  Let G be a group. One says that a \(\mathbf{Z}^{(G)}\) -module M is co-induced (resp. induced) if there exists a Z-module A such that M is isomorphic to \(\operatorname{Hom}_{\mathbf{Z}}(\mathbf{Z}^{(G)}, \mathbf{A})\) (resp. \(\mathbf{Z}^{(G)} \otimes_{\mathbf{Z}} \mathbf{A}\) ).
\end{enumerate}

\begin{enumerate}
\def\labelenumi{\alph{enumi})}
\item
  Every \(\mathbf{Z}^{(G)}\) -module is isomorphic to a submodule of a co-induced module and to a quotient of an induced module.
\item
  Let M be a \(\mathbf{Z}^{(G)}\) -module. Show that the following conditions are equivalent:
\end{enumerate}

α) M is a direct summand of a co-induced module.

β) M is injective relative to Z (exercise 19, p.~170).

γ) There exists a G-average on M (I, p.~136, exercise 8), that is, a \(\mathbf{Z}^{(G)}\) -homomorphism

\(m: \mathbf{C}^{1}(\mathbf{G}, \mathbf{M}) \to \mathbf{M}\) (the G-module structure on \(\mathbf{C}^{1}(\mathbf{G}, \mathbf{M})\) being defined by \((gf)(x) = gf(g^{-1}x)\) for \(g \in \mathbf{G}, x \in \mathbf{M}\) ) such that if \(f_{x}\) is the constant function with value \(x \in \mathbf{M}\) , one has \(m(f_{x}) = x\) .

\begin{enumerate}
\def\labelenumi{\alph{enumi})}
\setcounter{enumi}{2}
\tightlist
\item
  Show that the following conditions are equivalent:
\end{enumerate}

α) M is a direct summand of an induced module.

β) M is projective relative to Z (exercise 19, p.~170).

\(\gamma)\) There exists a \(\mathbf{Z}\) -endomorphism \(\rho\) of \(\mathbf{M}\) such that for every \(x \in \mathbf{M}\) , the family \((\rho(g^{-1} x))_{g \in G}\) has finite support and one has \(x = \sum g\rho(g^{-1} x)\) .

\begin{enumerate}
\def\labelenumi{\alph{enumi})}
\setcounter{enumi}{3}
\item
  Let M be an injective (resp. projective) \(\mathbf{Z}^{(G)}\) -module relative to \(\mathbf{Z}\) . Show that \(\mathrm{H}^q (\mathbf{G},\mathbf{M}) = 0\) (resp. \(\mathrm{H}_q(\mathrm{G},\mathrm{M}) = 0\) ) for every \(q\geqslant 1\) .
\item
  If G is finite, show that the notions of induced and co-induced module (resp. projective and injective relative to Z) coincide.
\end{enumerate}

\begin{enumerate}
\def\labelenumi{\arabic{enumi})}
\setcounter{enumi}{12}
\tightlist
\item
  Let G be a group, H a subgroup of G, M a left \(\mathbf{Z}^{(H)}\) -module, and N a right \(\mathbf{Z}^{(H)}\) -module. We endow the group \(\hat{\mathbf{M}} = \operatorname{Hom}_{\mathbf{Z}^{(H)}} (\mathbf{Z}^{(G)}, \mathbf{M})\) (resp. \(\check{\mathbf{N}} = \mathbf{N} \otimes_{\mathbf{Z}^{(H)}} \mathbf{Z}^{(G)})\) with the structure of \(\mathbf{Z}^{(G)}\) -module on the left (resp. on the right) deduced from the structure of \(\mathbf{Z}^{(G)}\) -module on the right of \(\mathbf{Z}^{(G)}\) .
\end{enumerate}

Let \(m:\hat{\mathbf{M}}\to \mathbf{M}\) and \(n:\mathbf{N}\rightarrow \check{\mathbf{N}}\) be the \(\mathbf{Z}^{(\mathrm{H})}\) -homomorphisms such that \(m(u) = u(1)\) for \(u\in \hat{\mathbf{M}}\) and

\[
n (x) = x \otimes 1 \quad \text { for } \quad x \in \mathbf {N}  ,
\]

and let \(i: H \to G\) be the canonical injection. Prove that the composed homomorphisms

\[
\hat {i} ^ {q}: \mathrm{H} ^ {q} (\mathrm{G}, \hat {\mathrm{M}}) \xrightarrow {i ^ {q}} \mathrm{H} ^ {q} (\mathrm{H}, \hat {\mathrm{M}}) \xrightarrow {\mathrm{H} ^ {q} (\mathrm{H} , m)} \mathrm{H} ^ {q} (\mathrm{H}, \mathrm{M})
\]

and

\[
\check {i} _ {q}: \mathrm{H} _ {q} (\mathrm{H}, \mathrm{N}) \xrightarrow {\mathrm{H} ^ {q} (\mathrm{H} , n)} \mathrm{H} _ {q} (\mathrm{H}, \check {\mathrm{N}}) \xrightarrow {i _ {q}} \mathrm{H} _ {q} (\mathrm{G}, \check {\mathrm{N}})
\]

are isomorphisms for all \(q \geqslant 0\) .

\begin{enumerate}
\def\labelenumi{\arabic{enumi})}
\setcounter{enumi}{13}
\tightlist
\item
  Let G be a group, and let H be a subgroup of G of finite index. We set \(n = (G:H)\) .
\end{enumerate}

\begin{enumerate}
\def\labelenumi{\alph{enumi})}
\tightlist
\item
  Let Q, M be two \(\mathbf{Z}^{(G)}\) -modules. Show that there exists a \(\mathbf{Z}\) -homomorphism
\end{enumerate}

\[
t: \operatorname{Hom} _ {\mathbf {Z} ^ {(H)}} (Q, M) \rightarrow \operatorname{Hom} _ {\mathbf {Z} ^ {(G)}} (Q, M)
\]

such that, for every system of representatives, \(g_{1}, \ldots, g_{n}\) of the classes of G modulo H,

\[
(t u) (q) = \sum_ {i = 1} ^ {n} g _ {i} u \left(g _ {i} ^ {- 1} q\right) \quad \text { for } \quad q \in \mathrm{Q}, \quad u \in \operatorname{Hom} _ {\mathbf {Z} (\mathrm{H})} (\mathrm{Q}, \mathrm{M}).
\]

If \(j: \operatorname{Hom}_{\mathbf{Z}(\mathbf{G})}(\mathbf{Q}, \mathbf{M}) \to \operatorname{Hom}_{\mathbf{Z}(\mathbf{H})}(\mathbf{Q}, \mathbf{M})\) is the canonical injection, we have \(t \circ j = n\) . Id.

\begin{enumerate}
\def\labelenumi{\alph{enumi})}
\setcounter{enumi}{1}
\tightlist
\item
  Let N be a right \(\mathbf{Z}^{(G)}\) -module. With the notation of \(a\) ), show that there exists a \(\mathbf{Z}\) -homomorphism \(\theta: \mathrm{N} \otimes_{\mathbf{Z}^{(G)}} \mathrm{Q} \to \mathrm{N} \otimes_{\mathbf{Z}^{(H)}} \mathrm{Q}\) such that
\end{enumerate}

\[
\theta (n \otimes q) = \sum_ {i = 1} ^ {n} n g _ {i} ^ {- 1} \otimes g _ {i} q \quad \text { for } \quad n \in \mathbf {N}, \quad q \in \mathbf {Q}.
\]

If \(p: \mathbf{N} \otimes_{\mathbf{Z}(\mathbf{H})} \mathbf{Q} \to \mathbf{N} \otimes_{\mathbf{Z}(\mathbf{G})} \mathbf{Q}\) is the canonical map, we have \(p \circ \theta = n\) . Id.
\begin{enumerate}
\def\labelenumi{\alph{enumi})}
\setcounter{enumi}{2}
\item
  Let P be a projective resolution of the \(\mathbf{Z}^{(G)}\) -module \(\mathbf{Z}\) , \(q\) an integer. By applying \(a\) and \(b)\) with \(\mathbf{Q} = \mathbf{P}_q\) , define homomorphisms \(t^q: \mathrm{H}^q(\mathrm{H}, \mathrm{M}) \to \mathrm{H}^q(\mathrm{G}, \mathrm{M})\) and \(\theta_q: \mathrm{H}_q(\mathrm{G}, \mathrm{N}) \to \mathrm{H}_q(\mathrm{H}, \mathrm{N})\) . If \(i\) denotes the canonical injection of H into G, we have \(t^q \circ i^q = n\) . Id. and \(i_q \circ \theta_q = n\) . Id.
\item
  We consider the \(\mathbf{Z}^{(G)}\) -modules \(\hat{\mathbf{M}} = \operatorname{Hom}_{\mathbf{Z}^{(H)}} (\mathbf{Z}^{(G)}, \mathbf{M})\) and \(\check{\mathbf{N}} = \mathbf{N} \otimes_{\mathbf{Z}^{(H)}} \mathbf{Z}^{(G)}\) (exercise 13). Taking \(\mathbf{Q} = \mathbf{Z}^{(G)}\) , we obtain from the homomorphisms \(t\) and \(\theta\) of \(a)\) and \(b)\) homomorphisms \(t_\mathbf{M}: \hat{\mathbf{M}} \to \mathbf{M}\) and \(\theta_\mathbf{N}: \mathbf{N} \to \check{\mathbf{N}}\) . Show that these homomorphisms are \(\mathbf{Z}^{(G)}\) -linear. Show that the composite homomorphism
\end{enumerate}

\[
\begin{array}{c} \mathrm{H} ^ {q} (\mathrm{H}, \mathrm{M}) \xrightarrow {\stackrel {{(\vec {l} _ {q}) ^ {- 1}}} {{\longrightarrow}}} \mathrm{H} ^ {q} (\mathrm{G}, \stackrel {{\hat {\mathrm{M}}}} {{\longrightarrow}}) \xrightarrow {\mathrm{H} ^ {q} (\mathrm{G} , t _ {\mathrm{M}})} \mathrm{H} ^ {q} (\mathrm{G}, \mathrm{M}) \\ (\text {resp.} \mathrm{H} _ {q} (\mathrm{G}, \mathrm{N}) \xrightarrow {\mathrm{H} _ {q} (\mathrm{G} , \theta_ {\mathrm{N}})} \mathrm{H} _ {q} (\mathrm{G}, \check {\mathrm{N}}) \xrightarrow {\stackrel {{(\vec {l} _ {q}) ^ {- 1}}} {{\longrightarrow}}} \mathrm{H} _ {q} (\mathrm{H}, \mathrm{N})) \text {is equal to } t ^ {q} (\text {resp.} \theta_ {q}). \end{array}
\]

\begin{enumerate}
\def\labelenumi{\alph{enumi})}
\setcounter{enumi}{4}
\item
  Let \((g_{1},\ldots ,g_{n})\) be a system of representatives of the classes of G modulo H, and let \(x\in \mathrm{H}^0 (\mathrm{H},\mathrm{M})\) . Show that we have \(t^0 (x) = \sum_{i = 1}^{n}g_{i}x\) . Describe \(t^1\)
\item
  If one identifies the groups \(\mathrm{H}_{1}(\mathrm{G},\mathbf{Z})\) and \(\mathrm{H}_{1}(\mathrm{H},\mathbf{Z})\) with G/(G, G) and H/(H, H) respectively, the homomorphism \(t_{1}\) is identified with a homomorphism (“transfer”) V : G/(G, G) → H/(H, H). If s : \(H \backslash G\) → G is a section of the canonical map, show that V is obtained by passage to the quotient from the map \(V_{s}\) : G → H such that
\end{enumerate}

\[
\mathrm{V} _ {s} (g) = \prod_ {x \in \mathrm{H} \backslash \mathrm{G}} s (x) g (s (x g)) ^ {- 1} \quad \text { for } \quad g \in \mathrm{G}.
\]

\begin{enumerate}
\def\labelenumi{\arabic{enumi})}
\setcounter{enumi}{14}
\tightlist
\item
  Let G be a group and k a commutative field, which we endow with the structures of \(\mathbf{Z}^{(\mathbf{G})}\) -module and of \(k^{(\mathbf{G})}\) -module deduced from the trivial action of G.
\end{enumerate}

\begin{enumerate}
\def\labelenumi{\alph{enumi})}
\item
  Show that the group \(\mathrm{H}^q (\mathbf{G},k)\) is isomorphic to \(\operatorname {Ext}_{k(\mathbf{G})}^{q}(k,k)\) for all \(q\geqslant 0\) , therefore admits a structure of \(k\) -vector space; we denote by \(h^q (\mathbf{G},k)\) its dimension over \(k\)
\item
  Suppose that the group G admits a presentation (1, p.~86) defined by d generators and r relations. Prove the inequalities
\end{enumerate}

\[
d \geqslant h ^ {1} (\mathbf {G}, k) \quad \text { and } \quad r \geqslant h ^ {2} (\mathbf {G}, k)
\]

(use exercise 3, d), p.~179).

¶ 16) Let G be a p-group (1, p.~72, def. 9). We set \(d(\mathbf{G}) = \dim_{\mathbf{F}_{p}} \mathrm{H}^{1}(\mathbf{G}, \mathbf{F}_{p})\) and \(r(\mathbf{G}) = \dim_{\mathbf{F}_{p}} \mathrm{H}^{2}(\mathbf{G}, \mathbf{F}_{p})\) (cf.~exercise 15).

\begin{enumerate}
\def\labelenumi{\alph{enumi})}
\item
  Show that \(d(\mathbf{G})\) is the minimal number of generators of \(\mathbf{G}\) (prove that \(\mathrm{H}^1 (\mathbf{G},\mathbf{F}_p)\) is isomorphic to the dual of \(\mathbf{G} / \mathbf{G}^p\mathbf{D}(\mathbf{G})\) , and use I, p.~140, exercise 32).
\item
  If Card (G) = \(p^n\) , show that one has \(d(\mathbf{G}) \leqslant n\) and \(r(\mathbf{G}) \leqslant \frac{1}{2} n(n + 1)\) (prove by induction on \(n\) that there exists a presentation of G with \(n\) generators and \(\frac{1}{2} n(n + 1)\) relations). If \(\mathbf{G} = (\mathbf{Z}/p\mathbf{Z})^n\) , show that one has \(d(\mathbf{G}) = n\) and \(r(\mathbf{G}) = \frac{1}{2} n(n + 1)\) .
\item
  If G is nontrivial, prove the inequality \(r(\mathbf{G}) > \frac{1}{4}(d(\mathbf{G}))^{2}\) (“Golod-Shafarevich theorem”: use exercise 15, p.~182, and exercise 22 of VIII, § 8).
\end{enumerate}

\(d)\) Prove that one has \(r(\mathbf{G}) - d(\mathbf{G}) = \dim_{\mathbf{F}_p}(\mathbf{H}^3 (\mathbf{G},\mathbf{Z})\otimes_{\mathbf{Z}}\mathbf{F}_p).\)

\begin{enumerate}
\def\labelenumi{\arabic{enumi})}
\setcounter{enumi}{16}
\tightlist
\item
  Let G be a finite group. We denote by N the element \(\sum_{g\in G}e_g\) of \(\mathbf{Z}^{(G)}\) . For every \(\mathbf{Z}^{(G)}\) -module M, the multiplication
\end{enumerate}

by N defines by passage to the quotient a homomorphism \(\mathrm{N}_{0}:\mathrm{H}_{0}(\mathrm{G},\mathrm{M})\to\mathrm{H}^{0}(\mathrm{G},\mathrm{M})\) . We set \(\hat{\mathrm{H}}^{0}(\mathrm{G},\mathrm{M})=\mathrm{Coker}(\mathrm{N}_{0}),\hat{\mathrm{H}}^{-1}(\mathrm{G},\mathrm{M})=\mathrm{Ker}(\mathrm{N}_{0}),\hat{\mathrm{H}}^{q}(\mathrm{G},\mathrm{M})=\mathrm{H}^{q}(\mathrm{G},\mathrm{M})\text{ for } q\geqslant1,\text{ and}\)

\[
\hat {\mathrm{H}} ^ {- q} (\mathrm{G}, \mathrm{M}) = \mathrm{H} _ {q - 1} (\mathrm{G}, \mathrm{M}) \quad \text { for } \quad q \geqslant 2.
\]

\begin{enumerate}
\def\labelenumi{\alph{enumi})}
\tightlist
\item
  Let \(\mathcal{E}:0\to\mathbf{M}^{\prime}\to\mathbf{M}\to\mathbf{M}^{\prime\prime}\to0\) an exact sequence of \(\mathbf{Z}^{(G)}\) -modules; define an exact sequence
\end{enumerate}

\[
\dots \rightarrow \hat {\mathrm{H}} ^ {q - 1} (\mathrm{G}, \mathrm{M} ^ {\prime \prime}) \xrightarrow {\partial^ {q - 1} (\mathrm{G} , \mathcal {E})} \hat {\mathrm{H}} ^ {q} (\mathrm{G}, \mathrm{M} ^ {\prime}) \rightarrow \hat {\mathrm{H}} ^ {q} (\mathrm{G}, \mathrm{M}) \rightarrow \hat {\mathrm{H}} ^ {q} (\mathrm{G}, \mathrm{M} ^ {\prime \prime}) \xrightarrow {\partial^ {q} (\mathrm{G} , \mathcal {E})} \hat {\mathrm{H}} ^ {q + 1} (\mathrm{G}, \mathrm{M} ^ {\prime}) \rightarrow \dots
\]

\begin{enumerate}
\def\labelenumi{\alph{enumi})}
\setcounter{enumi}{1}
\tightlist
\item
  If M is projective relative to Z, show that one has \(\hat{\mathrm{H}}^{q}(\mathrm{G},\mathrm{M}) = 0\) for all \(q\in \mathbf{Z}\) (cf.~exercise 12). c) Let H be a subgroup of G, \(i:\mathrm{H}\to \mathrm{G}\) the canonical injection. For every \(\mathbf{Z}^{(\mathrm{G})}\) -module M and every rational integer \(q\) one defines homomorphisms \(\hat{i}_{\mathrm{M}}^{\mathrm{q}}:\hat{\mathrm{H}}^{\mathrm{q}}(\mathrm{G},\mathrm{M})\rightarrow \hat{\mathrm{H}}^{\mathrm{q}}(\mathrm{H},\mathrm{M})\) and \(\hat{t}_{\mathrm{M}}^{\mathrm{q}}:\hat{\mathrm{H}}^{\mathrm{q}}(\mathrm{H},\mathrm{M})\rightarrow \hat{\mathrm{H}}^{\mathrm{q}}(\mathrm{G},\mathrm{M})\) (denoted simply \(\hat{i}^q\) and \(\hat{t}^q\) if there is no ambiguity about M) in the following way: for \(q\geqslant 1\) , we set \(\hat{i}^q = i^q\) and \(\hat{t}^q = t^q\) (exercises 10 and 14); for \(q\leqslant -2\) , we set \(\hat{i}^q = \theta_{-q - 1}\) and \(\hat{t}^q = i_{-q - 1}\) ; finally, one defines \(\hat{i}^0\) and \(\hat{t}^0\) (resp. \(\hat{i}^{-1}\) and \(\hat{t}^{-1}\) ) by passage to quotients (resp. to subsets) from \(i^0\) and \(t^0\) (resp. \(t_0\) and \(i_0\) ). Prove that one has \(\hat{i}^q\circ \hat{t}^q = n\) . Id. for all \(q\in \mathbf{Z}\) , where \(n\) is the index of H in G. Deduce that the groups \(\hat{\mathbf{H}}^{q}(\mathbf{G},\mathbf{M})\) are annihilated by the order of G (cf. exercise 6, p. 189).
\end{enumerate}

\(d)\) With the notations of \(b\) ), show that one has \(\hat{i}_{M'}^{q+1} \circ \partial^q(G, \mathcal{E}) = \partial^q(H, \mathcal{E}) \circ \hat{i}_{M''}^q\) and

\[
\hat {t} _ {\mathbf {M} ^ {\prime}} ^ {q + 1} \circ \partial^ {q} (\mathrm{H}, \mathcal {E}) = \partial^ {q} (\mathrm{G}, \mathcal {E}) \circ \hat {t} _ {\mathbf {M} ^ {\prime \prime}} ^ {q}
\]

for all \(q \in \mathbf{Z}\) .

\begin{enumerate}
\def\labelenumi{\alph{enumi})}
\setcounter{enumi}{4}
\tightlist
\item
  Let \(p\) be a prime number and H a \(p\) -Sylow subgroup (I, p. 74) of G. Show that the homomorphism \(\hat{i}^q: \hat{\mathrm{H}}^q(\mathrm{G}, \mathrm{M}) \to \hat{\mathrm{H}}^q(\mathrm{H}, \mathrm{M})\) has as kernel the sum of the \(l\) -primary components of the \(\mathbf{Z}\) -module \(\hat{\mathrm{H}}^q(\mathrm{G}, \mathrm{M})\) for \(l \neq p\) . In particular, if for every prime number \(p\) , one chooses a \(p\) -Sylow subgroup \(\mathrm{H}_p\) of G, the map \(\hat{\mathrm{H}}^q(\mathrm{G}, \mathrm{M}) \to \prod \hat{\mathrm{H}}^q(\mathrm{H}_p, \mathrm{M})\) deduced from the homomorphisms \(\hat{i}^q\) is injective.
\end{enumerate}

\begin{enumerate}
\def\labelenumi{\arabic{enumi})}
\setcounter{enumi}{17}
\tightlist
\item
  Let G be a finite group and M a \(\mathbf{Z}^{(G)}\) -module. One says that M is cohomologically trivial if for every subgroup H of G and for every \(q \in \mathbf{Z}\) , the group \(\hat{\mathrm{H}}^q(\mathrm{H}, \mathrm{M})\) is zero.
\end{enumerate}

\begin{enumerate}
\def\labelenumi{\alph{enumi})}
\item
  Show that a \(\mathbf{Z}^{(G)}\) -module projective relative to \(\mathbf{Z}\) (in particular an induced module, cf. exercise 12) is cohomologically trivial.
\item
  Let n, r, k be three integers such that r divides \((k^{n} - 1)\) ; let G be a cyclic group of order n, \(\sigma\) a generator of G. One endows the Z-module M = Z/rZ with the structure of \(\mathbf{Z}^{(G)}\) -module defined by \(\sigma m = km\) for all \(m \in M\) . Prove that M is projective relative to Z if and only if n and r are coprime.
\item
  With the notations of b), show that in order for \(\hat{\mathrm{H}}^{q}(\mathrm{G}, \mathrm{M}) = 0\) for all \(q \in \mathbf{Z}\) it is necessary and sufficient that one has \((r, k - 1)(r, 1 + k + \cdots + k^{n-1}) = r\) . In particular, if \(r = k^{n} - 1\) , show that M is cohomologically trivial.
\item
  From b) and c), deduce an example of a \(\mathbf{Z}^{(G)}\) -module cohomologically trivial which is not projective relative to \(\mathbf{Z}\) , and of a \(\mathbf{Z}^{(G)}\) -module M such that \(\hat{\mathrm{H}}^q (\mathrm{G},\mathrm{M}) = 0\) for all \(q\in \mathbf{Z}\) which is not cohomologically trivial.
\end{enumerate}

¶ 19) Let p be a prime number, G a p-group, M a Z \(^{(G)}\) -module. We set A = F \(_{p}^{(G)}\) .

\begin{enumerate}
\def\labelenumi{\alph{enumi})}
\tightlist
\item
  We assume that \(pM = 0\) . Show that the following conditions are equivalent:
\end{enumerate}

α) There exists a rational integer q such that \(\hat{\mathrm{H}}^{q}(\mathrm{G},\mathrm{M})=0\) .

β) M is cohomologically trivial (exercise 18).

\(\gamma)\) M is an induced \(\mathbf{Z}^{(G)}\) -module.

δ) The A-module M is free.

(To prove that \(\alpha\) ) implies \(\delta\) ), reduce to the case where \(q = -2\) ; show that one then has

\[
\operatorname{Tor} _ {1} ^ {\mathbf {A}} \left(\mathbf {F} _ {p}, \mathbf {M}\right) = 0,
\]

and conclude with exercise 4, p.~185 and exercise 22 of VIII, § 8).

\begin{enumerate}
\def\labelenumi{\alph{enumi})}
\setcounter{enumi}{1}
\tightlist
\item
  We assume that multiplication by \(p\) is injective in M. Show that the following conditions are equivalent:
\end{enumerate}

α) There exists an integer q such that the groups \(\hat{\mathrm{H}}^{q}(\mathrm{G},\mathrm{M})\) and \(\hat{\mathrm{H}}^{q+1}(\mathrm{G},\mathrm{M})\) are zero.

β) M is cohomologically trivial.

γ) The A-module M/pM is free.

Let G be a finite group, M a \(\mathbf{Z}^{(G)}\) -module.

\begin{enumerate}
\def\labelenumi{\alph{enumi})}
\tightlist
\item
  We assume that the \(\mathbf{Z}\) -module M is free. Show that the \(\mathbf{Z}^{(G)}\) -module M is projective if and only if for every Sylow subgroup H \(\subset\) G, the \(\mathbf{Z}^{(H)}\) -module M is cohomologically trivial.
\end{enumerate}

(To prove that the condition is sufficient, consider a resolution \(0 \to R \to L \to M \to 0\) of the \(\mathbf{Z}^{(G)}\) -module M, where L is free; show that the \(\mathbf{Z}^{(H)}\) -module \(Q = \operatorname{Hom}_{\mathbf{Z}}(M, N)\) is cohomologically trivial by observing that \(Q/pQ\) is cohomologically trivial and using Exercise 19. Deduce from Exercise 17, e) that we have \(H^1(G, Q) = 0\) and conclude from this that M is projective.)

\begin{enumerate}
\def\labelenumi{\alph{enumi})}
\setcounter{enumi}{1}
\tightlist
\item
  Show that the following conditions are equivalent:
\end{enumerate}

\(\alpha)\) The \(\mathbf{Z}^{(G)}\) -module M is cohomologically trivial.

β) The \(\mathbf{Z}^{(G)}\) module M admits a finite projective resolution.

γ) The \(\mathbf{Z}^{(G)}\) -module M admits a projective resolution of length one.

(To show that \(\alpha\) ) implies \(\gamma\) ), consider an exact sequence \(0 \to R \to L \to M \to 0\) , where \(L\) is a free \(\mathbf{Z}^{(G)}\) -module, and apply \(a)\) to \(R\) .)

\begin{enumerate}
\def\labelenumi{\alph{enumi})}
\setcounter{enumi}{2}
\item
  We assume that the Z-module M is divisible. Show that the \(\mathbf{Z}^{(G)}\) -module M is cohomologically trivial if and only if it is injective (consider an exact sequence \(0 \rightarrow M \rightarrow I \rightarrow Q \rightarrow 0\) , where I is an injective \(\mathbf{Z}^{(G)}\) -module, and show using b) that the \(\mathbf{Z}^{(G)}\) -module \(\operatorname{Hom}_{\mathbf{Z}}(\mathbf{Q}, \mathbf{M})\) is cohomologically trivial).
\item
  For a \(\mathbf{Z}^{(G)}\) -module to be cohomologically trivial, it is necessary and sufficient that it admit an injective resolution of length 1.
\end{enumerate}

¶ 21) Let G be a cyclic group of finite order.

\begin{enumerate}
\def\labelenumi{\alph{enumi})}
\item
  If M is a \(\mathbf{Z}^{(\mathbf{G})}\) -module and \(q\) an even (resp. odd) rational integer, the group \(\hat{\mathrm{H}}^q (\mathbf{G},\mathbf{M})\) is isomorphic to \(\hat{\mathrm{H}}^0 (\mathbf{G},\mathbf{M})\) (resp. \(\hat{\mathrm{H}}^1 (\mathbf{G},\mathbf{M}))\) (cf. exercise 7).
\item
  We denote by \(\mathcal{C}\) the set of classes of \(\mathbf{Z}^{(G)}\) -modules of finite type T such that the groups \(\hat{\mathrm{H}}^q (\mathbf{G},\mathbf{T})\) are finite for every \(q\in \mathbf{Z}\) ; if the \(\mathbf{Z}^{(G)}\) -module M is of type \(\mathcal{C}\) , we set
\end{enumerate}

\[
h (\mathrm{M}) = \operatorname{Card} \left(\mathrm{H} ^ {0} (\mathrm{G}, \mathrm{M})\right) / \operatorname{Card} \left(\mathrm{H} ^ {1} (\mathrm{G}, \mathrm{M})\right).
\]

Let \(0 \to M' \to M \to M'' \to 0\) be an exact sequence of \(\mathbf{Z}^{(G)}\) -modules. Show that if two of the modules \(M'\) , \(M, M''\) are of type \(\mathcal{C}\) , the same is true of the third, and that one then has \(h(M) = h(M') h(M'')\) . In particular, the function \(h\) defines a homomorphism from the Grothendieck group \(\mathbf{K}(\mathcal{C})\) to the multiplicative group \(\mathbf{Q}^*\) . c) Let \(M\) be a finite \(\mathbf{Z}^{(G)}\) -module. Show that \(M\) is of type \(\mathcal{C}\) and that one has \(h(M) = 1\) .

\begin{enumerate}
\def\labelenumi{\alph{enumi})}
\setcounter{enumi}{3}
\item
  Suppose that the order of G is a prime number p. We denote by \(\mathcal{C}'\) the set of classes of \(\mathbf{Z}^{(G)}\) -modules finitely generated in which multiplication by \(p\) has finite kernel and cokernel; if M is of type \(\mathcal{C}'\) , we set \(\varphi(M) = \text{Card (Coker } p_M)/\text{Card (Ker } p_M\) . Prove as in b) that \(\varphi\) defines a homomorphism from \(\mathbf{K}(\mathcal{C}')\) to \(\mathbf{Q}^*\) .
\item
  Show that \(\mathbf{K}(\mathcal{C}')\) is generated by the classes of \(\mathbf{Z}^{(G)}\) -modules \(T\) possessing one of the following properties:
\end{enumerate}

\begin{enumerate}
\def\labelenumi{(\roman{enumi})}
\item
  T is finite.
\item
  Multiplication by \(p\) in T is bijective.
\item
  T = Z endowed with the trivial action of G.
\item
  \(\mathbf{T} = \mathbf{Z}^{(\mathbf{G})}\) .
\item
  \(\mathbf{T} = \mathbf{Q}_p / \mathbf{Z}_p\) (cf. TG, III, p. 84, exercise 23), endowed with the trivial action of G.
\item
  \(\mathrm{T} = \mathbf{Z}^{(\mathrm{G})}\otimes_{\mathbf{Z}}(\mathbf{Q}_p / \mathbf{Z}_p)\) .
\end{enumerate}

(First show that every \(\mathbf{Z}^{(\mathrm{G})}\) -module of type \(\mathcal{C}'\) can be written in \(\mathbf{K}(\mathcal{C}')\) as a sum of a \(\mathbf{Z}^{(\mathrm{G})}\) -module finitely generated over \(\mathbf{Z}\) and a \(\mathbf{Z}^{(\mathrm{G})}\) -module in which multiplication by \(p\) is surjective. Then determine the structure of simple modules over the rings \(\mathbf{Q}^{(\mathrm{G})}\) and \(\mathbf{Q}_p^{(\mathrm{G})}\) .)

\begin{enumerate}
\def\labelenumi{\alph{enumi})}
\setcounter{enumi}{6}
\tightlist
\item
  Deduce from the preceding that one has \(\mathcal{C}' \subset \mathcal{C}\) and that if M is a \(\mathbf{Z}^{(G)}\) -module of type \(\mathcal{C}'\) , we have
\end{enumerate}

\[
h (\mathbf {M}) ^ {p - 1} = \left(\varphi \left(\mathrm{H} ^ {0} (\mathbf {G}, \mathbf {M})\right)\right) ^ {p} / \varphi (\mathbf {M}).
\]

(Verify the formula in the six cases considered in \(f\) ).)

\begin{enumerate}
\def\labelenumi{\arabic{enumi})}
\setcounter{enumi}{21}
\tightlist
\item
  Let K be a commutative field, L a finite Galois extension of K, G its Galois group. a) Show that the exact sequence of G-modules
\end{enumerate}

\[
0 \rightarrow \mu_ {n} (\mathrm{L}) \rightarrow \mathrm{L} ^ {*} \xrightarrow {x \mapsto x ^ {n}} \mathrm{L} ^ {* n} \rightarrow 0
\]

makes it possible to define an isomorphism \(k_{L}:(L^{n}\cap\mathbf{K}^{*})/\mathbf{K}^{*n}\to\mathbf{H}^{1}(\mathbf{G},\mu_{n}(\mathbf{L}))\) . When \(\mu_{n}(\mathbf{K})\) has \(n\) elements and one identifies \(\mathbf{H}^{1}(\mathbf{G},\mu_{n}(\mathbf{L}))\) with Hom \((\mathbf{G},\mu_{n}(\mathbf{K}))\) , show that \(k_{L}\) coincides with the homomorphism defined in V, § 11, n° 8.

\begin{enumerate}
\def\labelenumi{\alph{enumi})}
\setcounter{enumi}{1}
\tightlist
\item
  Let \(p\) be a prime number; suppose that \(K\) has characteristic \(p\) . Show that the exact sequence of G-modules (V, § 11, n° 9)
\end{enumerate}

\[
0 \to \mathbf {F} _ {p} \to \mathrm{L} \stackrel {{\mathfrak {p}}} {{\to}} \mathfrak {p} (\mathrm{L}) \to 0
\]

makes it possible to define an isomorphism

\[
a _ {L}: (\mathfrak {p} (L) \cap K) / \mathfrak {p} (K) \rightarrow H ^ {1} (G, F _ {p})
\]

which coincides with the homomorphism defined in loc. cit. if one identifies

\[
\mathbf {H} ^ {1} (\mathbf {G}, \mathbf {F} _ {p})
\]

with Hom

\[
(\mathbf {G}, \mathbf {F} _ {p})
\]

\begin{enumerate}
\def\labelenumi{\arabic{enumi})}
\setcounter{enumi}{22}
\tightlist
\item
  Let G be a group, M a group with operators in G (1, p. 29, def. 2). We denote by \(^{g}m\) the composite of \(g \in G\) and \(m \in M\) . We say that M is a G-group if one has \(^{h(g)m} = ^{hg}m\) for \(m \in M\) , \(g\) , \(h \in G\) , in other words if the map \(\varphi\) from G to Aut(M) given by \(\varphi(g), m = ^{g}m\) for \(g \in G\) , \(m \in M\) is a homomorphism.
\end{enumerate}

If M is a G-group, we denote by \(\mathbf{H}^{0}(\mathbf{G},\mathbf{M})\) the subgroup of M consisting of the elements \(m\in M\) such that \(^{g}m=m\) for all \(g\in G\) , and \(\mathbf{Z}^{1}(\mathbf{G},\mathbf{M})\) the set of maps h from G to M such that

\[
h (g g ^ {\prime}) = h (g) ^ {g} h (g ^ {\prime}) \quad \text { for } \quad g, g ^ {\prime} \in G.
\]

\begin{enumerate}
\def\labelenumi{\alph{enumi})}
\item
  We say that two elements \(h_1, h_2\) of \(Z^1(\mathbf{G}, \mathbf{M})\) are cohomologous if there exists an element \(m\) of \(\mathbf{M}\) such that \(h_2(g) = m^{-1} h_1(g)^g m\) for all \(g \in \mathbf{G}\) . Show that the relation thus defined is an equivalence relation on \(Z^1(\mathbf{G}, \mathbf{M})\) ; the quotient set is denoted by \(H^1(\mathbf{G}, \mathbf{M})\) . The class in \(H^1(\mathbf{G}, \mathbf{M})\) of the map \(i \in Z^1(\mathbf{G}, \mathbf{M})\) such that \(i(g) = 1_{\mathbf{M}}\) for all \(g \in \mathbf{G}\) is denoted by \(e_{G,\mathbf{M}}\) or simply \(e_{\mathbf{M}}\) .
\item
  Let \(f: \mathbf{M} \to \mathbf{N}\) be a homomorphism of G-groups (I, p.~30, def. 3); show that \(f\) induces a group homomorphism \(f^{0}: \mathrm{H}^{0}(\mathbf{G}, \mathbf{M}) \to \mathrm{H}^{0}(\mathbf{G}, \mathbf{N})\) , and a map \(f^{1}: \mathrm{H}^{1}(\mathbf{G}, \mathbf{M}) \to \mathrm{H}^{1}(\mathbf{G}, \mathbf{N})\) such that \(f^{1}(e_{\mathbf{M}}) = e_{\mathbf{N}}\) .
\item
  Let \(\mathbf{M} \xrightarrow{i} \mathbf{N} \xrightarrow{p} \mathbf{Q}\) be a sequence of \(\mathbf{G}\) -groups and homomorphisms, with \(i\) injective, \(p\) surjective and
\end{enumerate}

\[
\operatorname{Im} (i) = \operatorname{Ker} (p).
\]

Show that the group \(\mathbf{N}^{\prime}=p^{-1}(\mathbf{H}^{0}(\mathbf{G},\mathbf{Q}))\) acts on \(\mathbf{Z}^{1}(\mathbf{G},\mathbf{M})\) so that \((n,h)(g)=n^{-1}h(g)^{g}n\) for \(n\in\mathbf{N}^{\prime}\) , \(h\in\mathbf{Z}^{1}(\mathbf{G},\mathbf{M})\) , \(g\in\mathbf{G}\) ; show that this action induces, by passing to the quotient, an action of \(\mathbf{H}^{0}(\mathbf{G},\mathbf{Q})\) on \(\mathbf{H}^{1}(\mathbf{G},\mathbf{M})\) . The map \(i^{1}:\mathbf{H}^{1}(\mathbf{G},\mathbf{M})\to\mathbf{H}^{1}(\mathbf{G},\mathbf{N})\) passes to the quotient and defines a bijection from \(\mathbf{H}^{1}(\mathbf{G},\mathbf{M})/\mathbf{H}^{0}(\mathbf{G},\mathbf{Q})\) onto \(\operatorname{Im}(i^{1})\) .

\begin{enumerate}
\def\labelenumi{\alph{enumi})}
\setcounter{enumi}{3}
\tightlist
\item
  The action of \(\mathbf{H}^0 (\mathbf{G},\mathbf{Q})\) over \(\mathbf{H}^{1}(\mathbf{G},\mathbf{M})\) allows one to define a mapping \(\partial :\mathbf{H}^{0}(\mathbf{G},\mathbf{Q})\to \mathbf{H}^{1}(\mathbf{G},\mathbf{M})\) by \(\partial (q) = q\cdot e_{\mathrm{M}}\) for \(q\in \mathbf{H}^0 (\mathbf{G},\mathbf{Q})\) Show that we have
\end{enumerate}

\[
\operatorname{Im} \left(i ^ {0}\right) = \operatorname{Ker} \left(p ^ {0}\right), \operatorname{Im} \left(p ^ {0}\right) = \partial^ {- 1} (e _ {\mathrm{M}}), \quad \operatorname{Im} (\partial) = (i ^ {1}) ^ {- 1} \left(e _ {\mathrm{N}}\right), \quad \operatorname{Im} \left(i ^ {1}\right) = (p ^ {1}) ^ {- 1} \left(e _ {\mathrm{Q}}\right).
\]

\begin{enumerate}
\def\labelenumi{\alph{enumi})}
\setcounter{enumi}{4}
\item
  We now assume that \(i(\mathbf{M})\) is contained in the center of \(\mathbf{N}\) . Show that \(\partial\) is a group homomorphism; construct an action of the group \(\mathrm{H}^{1}(\mathrm{G},\mathrm{M})\) on the set \(\mathrm{H}^{1}(\mathrm{G},\mathrm{N})\) such that the map \(p^1\) passes to the quotient and defines a bijection from \(\mathrm{H}^{1}(\mathrm{G},\mathrm{N}) / \mathrm{H}^{1}(\mathrm{G},\mathrm{M})\) over \(\operatorname{Im}(p^1)\) .
\item
  Let E be the set of maps k from G to N such that \(p \circ k \in \mathbb{Z}^{1}(\mathbf{G}, \mathbf{Q})\) . For \(k \in E\) , show that there exists an element \(f_{k}\) of \(\mathbb{Z}^{2}(\mathbf{G}, \mathbf{M})\) such that \(i(f_{k}(x, y)) = k(x)^{x} k(y)(k(xy))^{-1}\) for \(x, y \in G\) . Show that the map \(k \mapsto f_{k}\) defines by passing to the quotient a map \(\partial^{1}: H^{1}(\mathbf{G}, \mathbf{Q}) \to H^{2}(\mathbf{G}, \mathbf{M})\) . Show that we have Im \((p^{1}) = (\partial^{1})^{-1}(0)\) .
\end{enumerate}

\begin{enumerate}
\def\labelenumi{\arabic{enumi})}
\setcounter{enumi}{23}
\tightlist
\item
  Let K be a commutative field, L a finite Galois extension of K, with Galois group G.
\end{enumerate}

Let V (resp. V') be a K-vector space, and t (resp. t') a tensor of type (p, q) on V (resp. V') (III, p.~63). A K-isomorphism (or simply isomorphism) of (V, t) onto (V', t') is a K-linear isomorphism u : V → V' such that (T\^{}p(u) ⊗ T\^{}q(ü))(t) = t'. We denote by V\_L the L-vector space L ⊗\_K V and t\_L is the tensor of type (p, q) on V\_L deduced from t by extension of scalars. We denote by F\_\{L/K\}(V, t) the set of isomorphism classes of pairs (W, τ) such that (W\_L, τ\_L) is L-isomorphic to (V\_L, t\_L).

\begin{enumerate}
\def\labelenumi{\alph{enumi})}
\tightlist
\item
  Let M be the group of L-automorphisms of \((V_{L}, t_{L})\) . For \(m \in M\) and \(\sigma \in G\) , we set
\end{enumerate}

\[
{ } ^ { \sigma } m = \pi _ { \mathrm{v} } ( \sigma ) \circ m \circ \pi _ { \mathrm{v} } ( \sigma ) ^ { - 1 } ,
\]

where \(\pi_{\mathbf{V}}: \mathbf{G} \to \operatorname{Aut}_{\mathbf{K}}(\mathbf{V}_{\mathbf{L}})\) denotes the action of \(\mathbf{G}\) on \(\mathbf{V}_{\mathbf{L}} = \mathbf{L} \otimes_{\mathbf{K}} \mathbf{V}\) deduced from the action of \(\mathbf{G}\) on \(\mathbf{L}\) . Show that one thus defines a structure of \(\mathbf{G}\) -group on \(\mathbf{M}\) .

\begin{enumerate}
\def\labelenumi{\alph{enumi})}
\setcounter{enumi}{1}
\item
  Let (W, τ) be an element of \(F_{L/K}(V, t)\) and let u be an L-isomorphism from \((\mathrm{V}_{\mathrm{L}}, t_{\mathrm{L}})\) onto \((\mathrm{W}_{\mathrm{L}}, \tau_{\mathrm{L}})\) . For \(\sigma \in G\) , we denote by \(f_{u}(\sigma)\) the element \(u^{-1} \circ \pi_{\mathbf{W}}(\sigma) \circ u \circ \pi_{\mathbf{V}}(\sigma)^{-1}\) of M. Show that \(f_{u} \in Z^{1}(\mathbf{G}, \mathbf{M})\) (Exercise 23), and that its class in \(H^{1}(\mathbf{G}, \mathbf{M})\) is independent of the choice of u; it is denoted by \(\theta(\mathbf{W}, \tau)\) .
\item
  Show that the map \(\theta : F_{L/K}(V, t) \to H^1(G, M)\) is bijective, and that one has \(\theta(V, t) = e_M\) .
\end{enumerate}

Let K be a commutative field, L a finite Galois extension of K, and G its Galois group. We let G act on the group GL(n, L) by setting σ(aij) = (σ(aij)) for σ ∈ G, (aij) ∈ GL(n, L). Every subgroup of GL(n, L) stable under this action is thereby endowed with a G-group structure (exercise 23). a) Show that the sets H¹(G, GL(n, L)) and H¹(G, SL(n, L)) are reduced to a single element (cf. V, § 10, no. 5, prop. 9).

\begin{enumerate}
\def\labelenumi{\alph{enumi})}
\setcounter{enumi}{1}
\item
  Show that the set \(\mathbf{H}^{1}(\mathbf{G},\mathbf{PGL}(n,\mathbf{L}))\) is identified with the set of isomorphism classes of central simple algebras S such that the algebra \(\mathbf{L}\otimes_{\mathbf{K}}\mathbf{S}\) is isomorphic to \(\mathbf{M}_{n}(\mathbf{L})\) (using Exercise 24). By associating to such an algebra its class in the group \(\mathbf{Br}(\mathbf{K},\mathbf{L})\) (VIII, § 13) and by identifying this latter group with \(\mathbf{H}^{2}(\mathbf{G},\mathbf{L}^{*})\) (loc. cit.), one obtains a map \(\mathbf{H}^{1}(\mathbf{G},\mathbf{PGL}(n,\mathbf{L}))\to \mathbf{H}^{2}(\mathbf{G},\mathbf{L}^{*})\) which is the opposite of the mapping \(\partial^{1}\) (Exercise 23, \(f\) ) deduced from the extension \(\mathbf{L}^{*}\to \mathbf{GL}(n,\mathbf{L})\to \mathbf{PGL}(n,\mathbf{L})\) . Show that the image of \(\partial^{1}\) is formed of elements of order \(n\) in \(\mathbf{H}^{2}(\mathbf{G},\mathbf{L}^{*})\) .
\item
  Show that the set \(\mathbf{H}^{1}(\mathbf{G},\mathbf{Sp}(2n,\mathbf{L}))\) is reduced to a single element (cf.~IX, § 5, no. 1, th. 1).
\item
  Let Q be a quadratic form on the K-vector space V, \(Q_{L}\) the quadratic form deduced from it on \(L \otimes_{K} V\) . Show that \(\mathrm{H}^{1}(\mathrm{G}, \mathbf{O}(\mathrm{Q}_{\mathrm{L}}))\) is identified with the set of isomorphism classes of quadratic forms \(Q'\) on V such that \(Q_{L}'\) is isomorphic to \(Q_{L}\) .
\end{enumerate}

\begin{enumerate}
\def\labelenumi{\arabic{enumi})}
\setcounter{enumi}{25}
\tightlist
\item
  We denote by \(\mathbf{A}^e\) the \(k\) -algebra \(\mathbf{A} \otimes_k \mathbf{A}^0\) ; every \((\mathbf{A}, \mathbf{A})\) -bimodule is naturally equipped with a structure of \(\mathbf{A}^e\) -module on the left and also a structure of \(\mathbf{A}^e\) -module on the right. Let M be a \((\mathbf{A}, \mathbf{A})\) -bimodule, \(n\) be an integer \(\geqslant 0\) ; we set:
\end{enumerate}

\[
\mathrm{H} _ {n} (\mathrm{A}, \mathrm{M}) = \operatorname{Tor} _ {n} ^ {\mathrm{A} ^ {e}} (\mathrm{M}, \mathrm{A}) \quad \mathrm{H} ^ {n} (\mathrm{A}, \mathrm{M}) = \operatorname{Ext} _ {\mathrm{A} ^ {e}} ^ {n} (\mathrm{A}, \mathrm{M}).
\]

\begin{enumerate}
\def\labelenumi{\alph{enumi})}
\item
  The \(k\) -module \(\mathbf{H}_0(\mathbf{A}, \mathbf{M})\) is isomorphic to the quotient of \(\mathbf{M}\) by the sub- \(k\) -module generated by the elements \((am - ma)\) for \(a \in \mathbf{A}, m \in \mathbf{M}\) ; the \(k\) -module \(\mathbf{H}^0(\mathbf{A}, \mathbf{M})\) is isomorphic to the sub- \(k\) -module of \(\mathbf{M}\) of elements \(m\) such that \(am = ma\) for all \(a \in \mathbf{A}\) . The \(k\) -module \(\mathbf{H}^1(\mathbf{A}, \mathbf{M})\) is identified with the quotient of the \(k\) -module of derivations \(\mathbf{D}_k(\mathbf{A}, \mathbf{M})\) by the sub- \(k\) -module of inner derivations, that is, the derivations \(d_m\) , for \(m \in \mathbf{M}\) , defined by \(d_m(a) = am - ma\) for all \(a \in \mathbf{A}\) (cf.~III, p.~132).
\item
  Assume that A is a projective k-module. Show that the standard resolution B(A) (X, p.~58) defines a projective resolution of the left \(\mathbf{A}^{e}\) -module A. Deduce that the k-modules H \(^{n}\) (A, M) are isomorphic to the homology modules of the complex C(A, M), where, for \(p \geqslant 0\) , C \(^{p}\) (A, M) is the k-module of k-multilinear maps from A \(^{p}\) to M, C \(^{p}\) (A, M) = 0 for \(p < 0\) , the differential being given by the formula :
\end{enumerate}

for \(f \in \mathbf{C}^p(\mathbf{A}, \mathbf{M})\) , \(x_0, \ldots, x_p \in \mathbf{A}\) .

Give an analogous expression for the k-modules \(H_{n}(A, M)\) .

\begin{enumerate}
\def\labelenumi{\alph{enumi})}
\setcounter{enumi}{2}
\item
  Under the hypothesis of \(b\) ), let \(\rho : k \to k'\) be a homomorphism of commutative rings, A' the \(k'\) -algebra \(k' \otimes_k \mathbf{A}\)', \(\mathbf{M}'\) a \((\mathbf{A}', \mathbf{A}')\) -bimodule. Show that the \(k\) -modules \(\mathrm{H}_n(\mathbf{A}, \mathbf{M}')\) and \(\mathrm{H}_n(\mathbf{A}', \mathbf{M}')\) (resp. \(\mathrm{H}^n(\mathbf{A}, \mathbf{M}')\) and \(\mathrm{H}^n(\mathbf{A}', \mathbf{M}')\) ) are isomorphic.
\item
  Assume henceforth that k is a field. Let B be a second k-algebra (associative and unital), M an (A, A)-bimodule, N a (B, B)-bimodule. For every \(n \geqslant 0\) , define an isomorphism of k-modules
\end{enumerate}

\[
\mathrm{H} _ {n} (\mathrm{A} \otimes_ {k} \mathrm{B}, \mathrm{M} \otimes_ {k} \mathrm{N}) \rightarrow \bigoplus_ {p + q = n} \left(\mathrm{H} _ {p} (\mathrm{A}, \mathrm{M}) \otimes_ {k} \mathrm{H} _ {q} (\mathrm{B}, \mathrm{N})\right).
\]

If one assumes moreover that A and B are finite-dimensional over k, show that there likewise exists a k-isomorphism

\[
\mathrm{H} ^ {n} (\mathrm{A} \otimes_ {k} \mathrm{B}, \mathrm{M} \otimes_ {k} \mathrm{N}) \rightarrow \bigoplus_ {p + q = n} \left(\mathrm{H} ^ {p} (\mathrm{A}, \mathrm{M}) \otimes_ {k} \mathrm{H} ^ {q} (\mathrm{B}, \mathrm{N})\right) \quad \text { for   every } \quad n \geqslant 0
\]

(use X, p.~76, th. 3).

\begin{enumerate}
\def\labelenumi{\alph{enumi})}
\setcounter{enumi}{4}
\tightlist
\item
  Let M, N be two left A-modules, P a right A-module, so that the k-modules \(\operatorname{Hom}_{k}(\mathbf{M}, \mathbf{N})\) and \(M \otimes_{k} P\) have natural structures of (A, A)-bimodules. Define for every \(n \geqslant 0\) isomorphisms of k-modules:
\end{enumerate}

\[
\mathrm{H} ^ {n} (\mathrm{A}, \operatorname{Hom} _ {k} (\mathrm{M}, \mathrm{N})) \rightarrow \operatorname{Ext} _ {\mathrm{A}} ^ {n} (\mathrm{M}, \mathrm{N}) \quad \text { and } \quad \mathrm{H} _ {n} (\mathrm{A}, \mathrm{M} \otimes_ {k} \mathrm{P}) \rightarrow \operatorname{Tor} _ {n} ^ {\mathrm{A}} (\mathrm{P}, \mathrm{M}).
\]

¶ 27) One calls an augmentation of the k-algebra A a unital homomorphism of k-algebras π : A → k; one says that the pair (A, π) is an augmented k-algebra. The augmentation endows k with a structure of left A-module. Let M be a left A-module, N a right A-module; one sets

\[
\mathrm{H} _ {n} (\mathrm{A}, \pi ; \mathrm{N}) = \operatorname{Tor} _ {n} ^ {\mathrm{A}} (\mathrm{N}, k) \quad \text { and } \quad \mathrm{H} ^ {n} (\mathrm{A}, \pi ; \mathrm{M}) = \operatorname{Ext} _ {\mathrm{A}} ^ {n} (k, \mathrm{M}) \quad \text { for   every } \quad n \geqslant 0.
\]

\begin{enumerate}
\def\labelenumi{\alph{enumi})}
\tightlist
\item
  We denote by \(\mathbf{M}_{(\pi)}\) (resp. \(\mathbf{N}_{(\pi)}\) ) the group \(\mathbf{M}\) (resp. \(\mathbf{N}\) ) endowed with the structure of (A, A)-bimodule defined by:
\end{enumerate}

\[
a m a ^ {\prime} = \pi (a ^ {\prime}) a m (\text { resp. } a n a ^ {\prime} = n \pi (a) a ^ {\prime} \quad \text { for } \quad a, \quad a ^ {\prime} \in \mathbf {A}, \quad m \in \mathbf {M}, \quad n \in \mathbf {N}).
\]

Show that the \(k\) -module \(\mathbf{H}_0(\mathbf{A}, \pi; \mathbf{N})\) (resp. \(\mathbf{H}^0(\mathbf{A}, \pi; \mathbf{M})\) , resp. \(\mathbf{H}^1(\mathbf{A}, \pi; \mathbf{M})\) ) is isomorphic to the \(k\) -module \(\mathbf{H}_0(\mathbf{A}, \mathbf{N}_{(\pi)})\) (resp. \(\mathbf{H}^0(\mathbf{A}, \mathbf{M}_{(\pi)})\) , resp. \(\mathbf{H}^1(\mathbf{A}, \mathbf{M}_{(\pi)})\) ) defined in exercise 26. \(b)\) We now assume that \(\mathbf{A}\) is a projective \(k\) -module. Define using the standard resolution (X, p.~58) the \(k\) -isomorphisms \(\varphi_n: \mathbf{H}_n(\mathbf{A}, \mathbf{N}_{(\pi)}) \to \mathbf{H}_n(\mathbf{A}, \pi; \mathbf{N})\) and \(\varphi^n: \mathbf{H}^n(\mathbf{A}, \mathbf{M}_{(\pi)}) \to \mathbf{H}^n(\mathbf{A}, \pi; \mathbf{M})\) for all \(n \geqslant 0\) .

\begin{enumerate}
\def\labelenumi{\alph{enumi})}
\setcounter{enumi}{2}
\item
  We assume that there exists a homomorphism of \(k\) -algebras \(\rho: \mathbf{A} \to \mathbf{A}^e\) such that \(\mathbf{A}^e\) is a projective right \(\mathbf{A}\) -module and such that the left ideal of \(\mathbf{A}^e\) generated by \(\rho(\text{Ker } \pi)\) be the kernel of the homomorphism \(\mathbf{A} \otimes_k \mathbf{A}^* \to \mathbf{A}\) defined by multiplication in \(\mathbf{A}\) . Show that for every \((\mathbf{A}, \mathbf{A})\) -bimodule \(\mathbf{Q}\) and every integer \(n \geqslant 0\) there are \(k\) -isomorphisms \(\mathrm{H}_n(\mathbf{A}, \pi; \mathbf{Q}) \to \mathrm{H}_n(\mathbf{A}, \mathbf{Q})\) (resp. \(\mathrm{H}^n(\mathbf{A}, \pi; \mathbf{Q}) \to \mathrm{H}^n(\mathbf{A}, \mathbf{Q})\) ), \(\mathbf{Q}\) being considered as \(\mathbf{A}\) -module on the right (resp. on the left) via the homomorphism \(\rho\) . (Use the canonical homomorphisms of \(\mathbf{X}\) , p.~109.)
\item
  Show that the hypotheses of c) are satisfied in the following cases:
\end{enumerate}

\(\alpha)\) A is the algebra (over \(k\) ) of a group \(\mathbf{G}\) , \(\rho\) is defined by \(\rho(e_g) = e_g \otimes e_{g-1}\) .

β) A is the tensor (resp. symmetric) algebra of a projective k-module V; ρ is defined by

\[
\rho (v) = v \otimes 1 - 1 \otimes v.
\]

* γ) A is the enveloping algebra of a Lie algebra g over k, free as a k-module; we have

\[
\rho (x) = x \otimes 1 - 1 \otimes x \quad \text { for } \quad x \in \mathfrak {g}
\]

(using cor. 5 of th. 1 of LIE I, § 2, no. 7).

δ) A is the ring of functions of a smooth algebraic group G over a field k; we take

\[
\rho (f) = \sum_ {i = 1} ^ {k} f _ {i} ^ {\prime} \otimes f _ {i} ^ {\prime \prime}, \quad \text { where } \quad f (g h ^ {- 1}) = \sum_ {i = 1} ^ {k} f _ {i} ^ {\prime} (g) f _ {i} ^ {\prime \prime} (h) \quad \text { for } \quad g, h \in G. *
\]

Let g be a Lie algebra over the commutative ring k, and U its enveloping algebra. We endow k with the structure of a left U-module corresponding to the trivial representation of g. Let N be a right g-module, M a left g-module, and n an integer ≥ 0; we set

\[
\mathrm{H} _ {n} (\mathfrak {g}, \mathrm{N}) = \operatorname{Tor} _ {n} ^ {\mathrm{U}} (\mathrm{N}, k) \quad \text { and } \quad \mathrm{H} ^ {n} (\mathfrak {g}, \mathrm{M}) = \operatorname{Ext} _ {\mathrm{U}} ^ {n} (k, \mathrm{M}).
\]

\begin{enumerate}
\def\labelenumi{\alph{enumi})}
\item
  The \(k\) -module \(\mathbf{H}_0(\mathfrak{g}, \mathbf{N})\) (resp. \(\mathbf{H}^0(\mathfrak{g}, \mathbf{M})\) ) is isomorphic to \(\mathbf{N}/\mathbf{Ng}\) (resp. the submodule of \(\mathbf{M}\) consisting of elements annihilated by \(\mathfrak{g}\) ). We denote by \(\mathbf{Z}^1(\mathfrak{g}, \mathbf{M})\) (resp. \(\mathbf{B}^1(\mathfrak{g}, \mathbf{M})\) ) the sub- \(k\) -module of \(\mathrm{Hom}_k(\mathfrak{g}, \mathbf{M})\) of elements \(f\) such that \(f([x, y]) = xf(y) - yf(x)\) for \(x, y \in \mathfrak{g}\) (resp. the elements \(f_m\) , for \(m \in \mathbf{M}\) , such that \(f_m(x) = xm\) for \(x \in \mathfrak{g}\) ). Show that \(\mathbf{H}^1(\mathfrak{g}, \mathbf{M})\) is identified with the quotient \(\mathbf{Z}^1(\mathfrak{g}, \mathbf{M})/\mathbf{B}^1(\mathfrak{g}, \mathbf{M})\) .
\item
  We now assume that \(\mathfrak{g}\) is a free \(k\) -module. Show that the \(k\) -modules \(\mathbf{H}^n(\mathfrak{g}, \mathbf{M})\) are identified with the homology modules of the complex \(\mathbf{C}(\mathfrak{g}, \mathbf{M})\) , where for \(p \geqslant 0\) , \(\mathbf{C}^p(\mathfrak{g}, \mathbf{M})\) is the \(k\) -module of alternating multilinear maps from \(\mathfrak{g}^p\) to \(\mathbf{M}\) , \(\mathbf{C}^p(\mathfrak{g}, \mathbf{M}) = 0\) for \(p < 0\) , the differential being given by the formula
\end{enumerate}

\[
\begin{array}{r l} d f (x _ {1}, \dots , x _ {p + 1}) = & \sum_ {1 \leqslant i \leqslant p + 1} (- 1) ^ {i + 1} x _ {i}. f (x _ {1}, \dots , \hat {x} _ {i}, \dots , x _ {p + 1}) \\ & + \sum_ {1 \leqslant i <   j \leqslant p + 1} (- 1) ^ {i + j} f ([ x _ {i}, x _ {j} ], x _ {1}, \dots , \hat {x} _ {i}, \dots , \hat {x} _ {j}, \dots , x _ {p + 1}) \end{array}
\]

for \(f \in C^{p}(\mathfrak{g}, \mathbf{M}), x_{1}, \ldots, x_{p+1} \in \mathfrak{g}\) (where the sign over a letter means that it is to be omitted).

Give an analogous expression for the \(k\) -modules \(\mathbf{H}_p(\mathfrak{g}, \mathbf{M})\) . (Use \(\mathbf{X}\) , p.~184, exercise 21.) c) We assume that \(\dim_k(\mathfrak{g}) = n < \infty\) . Show that \(\mathbf{H}^p(\mathfrak{g}, \mathbf{M}) = 0\) for all \(p \geqslant n + 1\) and for every \(\mathfrak{g}\) -module \(\mathbf{M}\) , and that there exists a \(\mathfrak{g}\) -module \(\mathbf{N}\) such that \(\mathbf{H}^n(\mathfrak{g}, \mathbf{N}) \neq 0\) (take \(\mathbf{N} = \Lambda^n\mathfrak{g}\) , where \(\mathfrak{g}\) acts by the formula

\[
x _ {*} \left(x _ {1} \wedge \dots \wedge x _ {n}\right) = \sum_ {i = 1} ^ {n} x _ {1} \wedge \dots \wedge x _ {i - 1} \wedge [ x, x _ {i} ] \wedge x _ {i + 1} \wedge \dots \wedge x _ {n};
\]

show that \(df = 0\) for all \(f \in \mathbf{C}^{n-1}(\mathfrak{g}, \mathbf{N})\) .

\begin{enumerate}
\def\labelenumi{\alph{enumi})}
\setcounter{enumi}{3}
\tightlist
\item
  The adjoint representation of g and the representation of g on M define a representation θ of g in C(g, M), such that
\end{enumerate}

\[
\left(\theta (x) f\right) \left(x _ {1}, \dots , x _ {p}\right) = x. f (x _ {1}, \dots , x _ {p}) - \sum_ {i = 1} ^ {p} f (x _ {1}, \dots , [ x, x _ {i} ], \dots , x _ {p})
\]

for \(f \in \mathbf{C}^p(\mathfrak{g}, \mathbf{M})\) , \(x, x_1, \ldots, x_p \in \mathfrak{g}\) . Show that, for every \(x \in \mathfrak{g}\) , \(\theta(x)\) is an endomorphism homotopic to zero of the complex \(\mathbf{C}(\mathfrak{g}, \mathbf{M})\) (define a homotopy \(i(x)\) by setting \((i(x) f) (x_1, \ldots, x_p) = f(x, x_1, \ldots, x_p)\) ). * 29) We keep the notation of the preceding exercise; we assume moreover that \(k\) is a field of characteristic zero, and that the Lie algebra \(\mathfrak{g}\) is semisimple (LIE, I, § 6). Let M be a g-module of finite dimension over \(k\) .

\begin{enumerate}
\def\labelenumi{\alph{enumi})}
\item
  If M is simple and nontrivial, then \(H^{p}(g, M)=0\) for every p\textgreater0 (the bilinear form on g associated with M is nondegenerate (LIE, I, § 6, n° 1, prop. 1); observe that multiplication by the corresponding Casimir element (LIE, I, § 3, n° 7) is zero in k and bijective in M).
\item
  Show that \(\mathbf{H}^{1}(\mathfrak{g},\mathbf{M}) = 0\) for every \(\mathbf{M}\) (use \(a\) ) and the fact that \(\mathfrak{g} = \mathcal{D}\mathfrak{g}\) ). Conversely, a \(k\) -Lie algebra \(\mathfrak{h}\) (of finite dimension) such that \(\mathbf{H}^{1}(\mathfrak{h},\mathbf{M}) = 0\) for all \(\mathfrak{h}\) -module \(\mathbf{M}\) of finite dimension is semisimple.
\item
  We give k the structure of g-module corresponding to the trivial representation. Show that, for every \(p \geqslant 0\) , the k-vector space \(H^{p}(g, k)\) is isomorphic to the subspace of \(C^{p}(g, k)\) consisting of alternating p-linear forms f that are invariant, that is, such that
\end{enumerate}

\[
\sum_ {i = 1} ^ {p} f (x _ {1}, \dots , [ x, x _ {i} ], \dots , x _ {p}) = 0 \quad \text { for   all } \quad x, x _ {1}, \dots , x _ {p} \in \mathfrak {g}.
\]

(Deduce from the semisimplicity of the representations of g and from exercise 28, d) that \(\mathbf{H}^p(\mathfrak{g}, k)\) is isomorphic to \(\mathbf{H}^p(\mathbf{C}_{\theta}(\mathfrak{g}, k))\) , where \(\mathbf{C}_{\theta}(\mathfrak{g}, k)\) is the subcomplex of \(\mathbf{C}(\mathfrak{g}, k)\) annihilated by all the endomorphisms \(\theta(x)\) for \(x \in \mathfrak{g}\) ; then show that \(\mathbf{C}_{\theta}(\mathfrak{g}, k)\) has zero differential.)

\begin{enumerate}
\def\labelenumi{\alph{enumi})}
\setcounter{enumi}{3}
\tightlist
\item
  Deduce from c) that \(\mathrm{H}^{2}(\mathrm{g},\mathrm{M})=0\) for every g-module M of finite dimension over k. (Show that every invariant bilinear form on g is symmetric, by reducing to the case where g is simple and k is algebraically closed.) *
\end{enumerate}
